\documentclass[11pt,a4paper]{article}
\usepackage[T1]{fontenc}
\usepackage{lmodern,amsmath,amssymb,amsthm,mathtools,booktabs,array,microtype,xcolor}
\usepackage[margin=26mm]{geometry}
\usepackage{needspace}
\usepackage{hyperref}
\hypersetup{hypertexnames=false}
\hypersetup{colorlinks=true,linkcolor=blue!50!black,citecolor=blue!50!black,urlcolor=blue!50!black,pdftitle={Fractional evolution, gamma quotients, and normalized Taylor remainders},pdfauthor={Research draft prepared with ChatGPT for Keio}}
\newtheorem{theorem}{Theorem}[section]
\newtheorem{lemma}[theorem]{Lemma}
\newtheorem{corollary}[theorem]{Corollary}
\newtheorem{proposition}[theorem]{Proposition}
\numberwithin{equation}{section}
\setlength{\emergencystretch}{2em}



\title{Fractional evolution, gamma quotients,\protect\\and normalized Taylor remainders}
\author{Research manuscript prepared with ChatGPT for Keio}
\date{3 November 2026}
\begin{document}
\maketitle
\begin{abstract}
Initial states for a fractional Schr\"odinger equation can give identical
terminal observations at exceptional orders and times. We construct two
infinite families of such exceptions, one for a scalar multiplier and
one for the determinant governing two initial data. Each has a separate
finite counterexample certified by interval arithmetic. At orders at most
one half a Gaussian scale-mixture representation gives logarithmic
interpolation of the solution norm. For a different fractional equation
we certify the order at which the first positive zero appears, its
square-root splitting, and its drift at the integer endpoint.

Two gamma-quotient questions lead to a comparison of Laplace densities
and a sharp threshold for a reciprocal. For normalized Taylor remainders
we prove sign theorems for squared tangent and arctangent, determine
optimal constants for products indexed by majorized vectors, and show
that raising the arcsine quotient to a large power destroys logarithmic
absolute monotonicity at every fixed truncation. A final graph theorem
gives a sharp linear leaf bound. Appendices treat typical shifts in
Erd\H{o}s--Graham Problem~461 and classify periodic free-isospectral
potentials on at most four occupied sites.
\end{abstract}
\medskip
\noindent\textbf{How to read the claims.}
The ten numbered entries on the next page are conjectural statements
posed in the cited sources. We give mathematical proofs or disproofs
in their original ranges. They share techniques, and their ages and
difficulty vary. The number counts original statements, not distinct
historically unprecedented methods. The source comparisons below
identify known ingredients, alternative proofs, and specific priority
uncertainties. External mathematical review and historical verification
remain open.
\clearpage
\section{Ten source conjectures and their outcomes}
The page references lead to the proofs. In entries 4--5, LCM means
logarithmic complete monotonicity. Two separately published questions
about reciprocal tangent quotients, entries 6--7, share one proof.
Entry 8 includes three related inequalities in the authors' Fifth
Conjecture and is counted once.
\smallskip
\begin{center}\small
\begin{tabular}{@{}r>{\raggedright\arraybackslash}p{.315\textwidth}>{\raggedright\arraybackslash}p{.518\textwidth}r@{}}
\toprule
& Original statement & Result and scope & Page\\
\midrule
1 & Chorfi et al., Conj.~2 (2025) & False. Infinitely many singular terminal-data determinants, with order asymptotics and a certified finite example. & \pageref{sec:wronskian}\\[5pt]
2 & Chorfi et al., Conj.~3 (2025) & True. A probability mixture proves complete monotonicity for every $0<a\le1/2$ and the proposed norm interpolation. & \pageref{sec:schrod3}\\[5pt]
3 & Chaudhary et al., Conj.~2 (2024) & False as written. The endpoint coefficient exceeds $1.01$. The threshold, strict monotonicity and square-root law hold. & \pageref{sec:drift}\\[5pt]
4 & Guo--Qi, Remark~4 (2011) & True. The normalized gamma-ratio derivative or its reciprocal is LCM in the respective full shift ranges. & \pageref{sec:gqderiv}\\[5pt]
5 & Qi--Wei--Guo, Remark~4.5 (2012) & False. An explicit counterexample and the sharp reciprocal LCM threshold for every shift. & \pageref{sec:qwg}\\[5pt]
6 & Liu--Qi, Remark~4.1 (2025) & True. Strict concavity of the higher divided by the lower squared-tangent remainder on $(-\pi/2,\pi/2)$. & \pageref{sec:tan-tails}\\[5pt]
7 & Shuang--He--Qi, Guess~2 (2025) & True. All derivative signs of the reciprocal tangent quotient on both half-intervals. & \pageref{sec:tan-tails}\\[5pt]
8 & Qi, Fifth Conjecture (2024--26 sources) & True. A comparison on constant-sum beta slices also gives both sharp constants in He--Qi Problem~3. & \pageref{sec:hqmajor}\\[5pt]
9 & Qi, Sixth Conjecture (2025) & False. Exact counterexample at $q=56$ and opposite eventual behavior in the two parameter directions. & \pageref{qi6:sec}\\[5pt]
10 & Liu--Qi, Seventh Conjecture in Qi's survey (2026) & True. Strict log-concavity of every normalized arctangent tail. The common upper endpoint $1$ is sharp. & \pageref{sec:qi-seventh}\\
\bottomrule
\end{tabular}\end{center}
\clearpage
\noindent\textbf{Further exact conclusions and their status.}
Chorfi et al.'s Conjecture~1 is false. Section~\ref{sec:imaginary}
gives infinitely many imaginary-axis crossings and a certified
enclosure. Hilfer and Seybold had already plotted zero trajectories
indicating the crossings in 2006. For the 2005 Graffiti.pc
Conjecture~165, Section~\ref{wow165:sec} gives a sharp spanning-tree
bound. A 2006 conference abstract by Foster prevents a firm priority
judgment. Section~\ref{sec:qiluo} solves a two-parameter open
classification question of Qi and Luo.

Chaudhary et al.'s Conjectures~1 and~3 contain exact inequalities and
numerical fits with no specified tolerance. Section~\ref{sec:cdh13}
proves the inequalities and gives exact asymptotic information in place
of the fits. Section~\ref{sec:normalized-remainder-oscillation}
refutes both parts of Shuang--He--Qi's joint Guess~3. Together with
Section~\ref{sec:moment-ratios} the remainder sections cover the
nine statements listed in Section~12.7.1 of Qi's survey
\href{https://arxiv.org/abs/2512.15797}{arXiv:2512.15797}.
Two are accessible from older theorems, as explained in
Section~\ref{sec:moment-ratios}.

\medskip
\noindent\textbf{Corpus and proof design.}
We consulted the \href{https://github.com/QEEDDD/Verifiable-Mathematical-Reasoning-Corpus/commit/ecd87de3348dcb706a76bb12fde722d7a348a012}{Verifiable Mathematical Reasoning Corpus,
revision \texttt{ecd87de3348dc}}. The restricted leaf estimate in [066]
suggests asking how a tree assembled from components pays for neighbors
of maximum-degree vertices, which yields the general graph result.
The proof-transfer examples in [067] motivate separate auxiliary
questions and explicit maps back to the original variables.
The sparse Floquet studies [017,061] enter Appendix~\ref{sec:four-sites}.
Each main proof specifies its own analytical or combinatorial inputs.
The earlier ten-entry manuscript is separate.

\medskip
\noindent\textbf{Two transfers.}
For low fractional orders, the auxiliary problem is to describe the
difference of two independent positive half-normal scale mixtures.
Its solution produces the probability mixture for an evolution
multiplier, then a norm bound for its Hilbert-space equation. For
orders near two, a decaying second residue is invisible in a fixed-order
asymptotic. Letting order and time vary independently balances that
residue with the algebraic part of a determinant. A winding argument
then proves loss of uniqueness for the original two initial data.
\medskip

\noindent\textbf{Computation.}
Four executable interval verifiers are included after the end of the
typeset material in the single TeX file. The two nonvanishing conjectures have independent analytic disproofs and finite interval certificates. The fractional positivity and fold bounds use
their indicated analytic estimates and rigorous interval calculations.
\clearpage


\section{Imaginary-axis crossings near the exponential order}
\label{sec:imaginary}
Chorfi, Et-tahri, Maniar and Yamamoto, \emph{Forward and backward
problems for abstract time-fractional Schr\"odinger equations},
\href{https://doi.org/10.1016/j.cnsns.2026.109758}{Communications
in Nonlinear Science and Numerical Simulation 157 (2026), 109758},
\href{https://arxiv.org/abs/2510.03600}{arXiv:2510.03600},
Section~5, Conjecture~1, ask whether $E_{a,1}(ix)$ is nonzero for
every real $x$ and every $3/5<a<1$. Here $E_{a,b}(z)=
\sum_{j\ge0}z^j/\Gamma(aj+b)$. A phase variable and an independent
damping variable give zeros approaching the order $1$.

\begin{theorem}\label{th:imaginary}
For every sufficiently large positive integer $k$ there are
$3/5<a_k<1$ and $x_k>0$ with $E_{a_k,1}(ix_k)=0$.
They may be chosen so that $x_k\sim 2\pi k$ and
\[
 \pi^2 k(1-a_k)=2\log k-\log\log k+3\log\pi+o(1).
\]
In particular infinitely many distinct orders give counterexamples.
\end{theorem}
\begin{proof}
Let $a=1-d$, $0<d\le1/10$, and $z=iX$ with $X\ge1$.
Deformation of the inverse Laplace contour for $E_{a,1}(zu^a)$
at $u=1$ gives
\[
 E_{a,1}(z)=a^{-1}e^{z^{1/a}}+J_a(z),\qquad
 J_a(z)=-\frac{z\sin(\pi a)}{\pi}
 \int_0^\infty\frac{e^{-r}r^{a-1}}
 {(r^ae^{i\pi a}-z)(r^ae^{-i\pi a}-z)}\,dr .
\]
Indeed the transform $s^{a-1}/(s^a-z)$ on the principal sheet has
one simple pole at $z^{1/a}$ of residue $1/a$. The two sides of the
negative real cut give the displayed difference. An inner arc has
integral $O(\epsilon^a)$, and on a left semicircle of large radius
$R$ the transform is $O(R^{-1})$. The factor $e^s$ makes its arc
integral tend to zero by the usual Jordan estimate. One may start
on a line to the right of all singularities and close about its
intersection with the real axis.

The angular separations imply
$|r^ae^{\pm i\pi a}-iX|\ge c(r^a+X)$ for an absolute $c>0$.
The reciprocal of the product differs from $1/z^2$ by at most
$C r^a/X^3$. This follows by subtracting the fractions and bounding
the resulting numerator by $2Xr^a+r^{2a}$. Integration and Euler's
reflection formula now give, uniformly in this range,
\[
 J_a(iX)=\frac{i}{X\Gamma(1-a)}+O(d/X^2)
 =\frac{i}{X\Gamma(d)}(1+O(X^{-1})).
\]
In particular the algebraic term stays in a fixed sector around the
positive imaginary axis, even when $d$ approaches zero.

Put $\phi=\pi/(2a)$, use coordinates $(d,\theta)$ and let
$t=\theta/\sin\phi$, $X=t^a$. On the rectangle
\[
 k^{-2}\le d\le\log(k)/k,\qquad
 2\pi k+\pi\le\theta\le2\pi k+2\pi ,
\]
one has $t\sim2\pi k$ and $X\sim2\pi k$ uniformly. Define the real
continuous functions
\[
 U=t\cos\phi-\log(a|J_a(iX)|),\qquad
 V=\theta-2\pi k-\arg[-J_a(iX)],
\]
taking the branch of argument near $3\pi/2$. On the two $\theta$
faces $V$ has opposite strict signs. On the lower $d$ face,
$t\cos\phi=O(k^{-1})$ and $a|J_a(iX)|\asymp k^{-3}$, so $U>0$.
On the upper face,
$t\cos\phi=-(\pi^2+o(1))\log k$ and
$\log(a|J_a(iX)|)=-2\log k+O(\log\log k)$, so $U<0$.
The two-dimensional intermediate value theorem (Poincar\'e--Miranda)
gives a common zero. At it
$e^{(iX)^{1/a}}=-aJ_a(iX)$, proving the required cancellation.

Throughout this rectangle
$t\cos\phi=-\pi^2kd+o(1)$ because $kd^2=o(1)$, and
$\log(a|J_a(iX)|)=\log d-\log(2\pi k)+o(1)$.
At a common zero $w=\pi^2kd$ therefore satisfies
$w+\log w=2\log k+\log(2\pi^3)+o(1)$.
First $w\sim2\log k$, then substitution gives the stated expansion.
Since $d>0$ tends to zero, the orders cannot form a finite set.
\end{proof}

On a scalar eigenmode with eigenvalue $\lambda>0$, the time-fractional
equation $i\partial_t^a u-\lambda u=0$ (Caputo derivative) has
$u(T)=E_{a,1}(-i\lambda T^a)u(0)$, as follows by a Laplace transform.
By complex conjugation, for each crossing above, $T=(x_k/\lambda)^{1/a_k}$ gives identical
terminal data for every initial amplitude. Thus an all-time uniqueness
claim for the backward problem already fails on a single eigenmode.

\paragraph{Earlier zero trajectories.}
R. Hilfer and H. J. Seybold, \emph{Computation of the generalized
Mittag--Leffler function and its inverse in the complex plane},
\href{https://doi.org/10.1080/10652460600725341}{Integral Transforms
and Special Functions 17 (2006), 637--652}, Figure~1 and
Theorems~3.1--3.2, plotted zero trajectories across the orders and
studied endpoint limits for fixed-index zeros. Their plotted paths
already indicate imaginary-axis crossings. This prior work is
material to the 2025 conjecture. The preceding direct proof gives
a large-index law for the crossing orders and the resulting
nonuniqueness times. Historical priority for the qualitative
observation is not asserted.


\paragraph{A finite counterexample enclosure.}
For an independent verification of the zero, set $F(a,t)=(\Re E_{a,1}(it),\Im E_{a,1}(it))$,
\[
 \begin{gathered}
 x_0=(0.7478564279599293430791082426,\,
       3.7131384963793705236713833185),\\
 B=10^{-3}\begin{pmatrix}424&-523\\8133&792\end{pmatrix}.
 \end{gathered}
\]
On $D=x_0+[-10^{-12},10^{-12}]^2$, outward-rounded interval evaluation gives
\[
 \|BF(x_0)\|_\infty<3\cdot10^{-29},\qquad
 \sup_{x\in D}\|I-BDF(x)\|_\infty<0.000132.
\]
Thus $x\mapsto x-BF(x)$ contracts $D$ into itself. Since $B$ is invertible, its fixed point gives a zero in $D$.

These two bounds can be reproduced by evaluating the first $240$ terms of the complex series at 160 bits of Arb precision, multiplying each summand $v_k=(it)^k/\Gamma(1+ak)$ by $-k\psi(1+ak)$ and $k/t$ for the two derivative columns, respectively. For $k\ge240$ each resulting summand is bounded in modulus by $(k+1)^2 4^k/\lfloor k/2\rfloor!$, since $1/2<a<1$, $1<t<4$ and $0<\psi(1+ak)<\log(1+ak)$. For $k=2m,2m+1$ together the bound is $5(2m+2)^2 4^{2m}/m!$. Consecutive such bounds have ratio below $1/4$ for $m\ge120$. Thus a radius $10^{-42}$ for each real or imaginary tail suffices.


\section{Zeros of the backward evolution determinant}
\label{sec:wronskian}
Conjecture~2 in Section~5 of Chorfi et al. concerns the determinant
in their equation (4.14). With $\lambda=i^{-a}=e^{-i\pi a/2}$, it is
\[
 \Psi_a(t)=E_{a,1}(\lambda t)^2
 -\lambda t E_{a,2}(\lambda t) E_{a,a}(\lambda t).
\]
They conjecture $\Psi_a(t)\ne0$ for every $t>0$ and $1<a<2$.
The one-pole expansion for fixed order in their Lemma~4.10 leaves
out an exponentially small contribution. As $a$ approaches $2$ that
contribution can cancel the algebraic term.
\begin{theorem}\label{th:wronskian}
Let $N_k=2\pi k+3\pi/2$. For all sufficiently large positive integers
$k$ there exist $a_k\in(1,2)$ and $w_k=N_k+o(1)$ with
\[
 \Psi_{a_k}(w_k^{a_k})=0,\qquad
 \pi^2 k(2-a_k)=2\log k-\log\log k+\log(2\pi^3)+o(1).
\]
\end{theorem}
\begin{proof}
Write $d=2-a$, $y_b(w)=w^{b-1}E_{a,b}(\lambda w^a)$ for $b=1,2$.
Termwise differentiation gives $y_2'=y_1$ and
$y_1'=\lambda w^{a-1}E_{a,a}(\lambda w^a)$. Thus the original
determinant is $y_1^2-y_1'y_2$ at $t=w^a$.

On $1.9\le a<2$ the two simple principal-sheet poles of the transforms
$s^{a-b}/(s^a-\lambda)$ are
\[
 \xi_0=-i,\qquad \xi_1=-\sin\theta+i\cos\theta,qquad
 \theta=\pi d/a.
\]
Set $e_j=e^{\xi_j w}$. For $w\ge1$ one has uniformly
\begin{align*}
 y_1&=(e_0+e_1)/a+O(w^{-a}),\\
 y_1'&=(\xi_0 e_0+\xi_1 e_1)/a+O(w^{-a-1}),\\
 y_2&=(e_0/\xi_0+e_1/\xi_1)/a-
          \frac{w^{1-a}}{\lambda\Gamma(2-a)}+O(w^{1-2a}).
\end{align*}
For the bounds deform the inverse Laplace contour about the negative
real cut and take the two residues. On its two sides $s=re^{\pm i\pi}$
one has $|s^a-\lambda|\ge c(1+r^a)$ with an absolute constant $c>0$
by angular separation. The $y_1$ cut is bounded by a constant times
$\int_0^\infty e^{-wr}r^{a-1}\,dr=O(w^{-a})$.
Differentiation there inserts one factor $-r$ and proves the derivative
estimate. For $y_2$ subtract $-s^{a-2}/\lambda$ from the cut
integrand. The remaining integrand before multiplication by $e^{sw}$
is bounded by $C r^{2a-2}/(1+r^a)$.
The inverse cut of the subtracted power is
$-w^{1-a}/(\lambda\Gamma(2-a))$ by the gamma integral.
The inner arcs vanish as $a-2>-1$. The outer arcs vanish by
exponential decay away from the ends and the $O(1/|s|)$ or stronger
decay of the original transforms. All constants are bounded uniformly
over the indicated orders.

The individual squared exponentials in $y_1^2-y_1'y_2$ cancel.
Define
\[
 C_a=\frac{2-\xi_0/\xi_1-\xi_1/\xi_0}{a^2}
     =\frac{2+2\cos\theta}{a^2},\qquad
 B_a=\frac{\xi_0}{a\lambda\Gamma(d)}.
\]
As $d\downarrow0$, $C_a\to1$ and $B_a\sim i d/2$.
Whenever $w\to\infty$, $d\to0$, $dw\to\infty$ and
$|e_1|=O(d w^{1-a})$, the preceding estimates give
\[
 e^{iw}\Psi_a(w^a)=C_a e_1+B_a w^{1-a}+o(d w^{1-a}).
\]
Indeed, the total remaining error is
$O(w^{-a}+w^{1-2a}+|e_1|d w^{1-a})$. These bounds also give
uniform little-oh estimates on the windows used below.

For $N=N_k$ choose $d_N$ between $2\log N/(\pi N)$ and
$6\log N/(\pi N)$ so that, with $a_N=2-d_N$,
\[
 C_{a_N}\exp[-N\sin(\pi d_N/a_N)]
     = |B_{a_N}|N^{1-a_N}=:D_N.
\]
Such a choice follows by continuity. At the two endpoints the left
side has orders $N^{-1}(1+o(1))$ and $N^{-3}(1+o(1))$, while
the right side has order $(\log N)/N^2$ at both. Taking logarithms
in the equality gives $d_N\sim4\log N/(\pi N)$.

Use real coordinates $d=d_N-2s/(\pi N)$, $w=N+v$. On any fixed
bounded disk in $(s,v)$ the exponential modulus in the preceding
identity is multiplied by $e^{s+o(1)}$, since
$[\sin(\pi d/(2-d))]'\to\pi/2$. The phase
$w\cos\theta=N+v+o(1)$ because $Nd^2=o(1)$.
Also $B_a w^{1-a}/D_N\to i$ uniformly, since $d/d_N\to1$ and
$(a-a_N)\log N=o(1)$. Consequently the continuous maps satisfy
\[
 D_N^{-1}e^{iw}\Psi_a(w^a)\longrightarrow
 i(1-e^{s+iv})
\]
uniformly on this disk. On every sufficiently small circle about the
origin the limit has winding number one, by its simple zero there.
Uniform convergence preserves the winding number and forces an interior
zero of the exact real two-parameter map. Choose a zero on a fixed
closed disk of radius $1/4$ for each large $k$. Every limit point of
these coordinates is a zero of the limit map on that disk and hence
is the origin. This proves $w_k=N_k+o(1)$. At these zeros the
modulus equation gives
$\exp[-(\pi/2)Nd+o(1)]=(1+o(1))d/(2N)$ because $Nd^2=o(1)$.
Thus $Q=(\pi/2)Nd$ satisfies $Q+\log Q=2\log N+\log\pi+o(1)$.
Using $Q\sim2\log N$ and $N=2\pi k+O(1)$ gives the claimed expansion.
\end{proof}

For an eigenvalue of magnitude one the scalar map from the two initial
data to the terminal value and terminal derivative has determinant
$\Psi_a(w^a)$ by the representation in the source. At the zeros in
Theorem~\ref{th:wronskian} it has a nontrivial kernel. Together with
Theorem~\ref{th:imaginary}, this illustrates two different failures of
terminal uniqueness as the order and observation time vary. For any fixed
order the source already proves finiteness of the exceptional positive
arguments (Lemma~4.11). The infinitude here ranges across orders.


\paragraph{A certified finite counterexample.}
In the coordinates $(a,w)$ with original argument $t=w^a$, the
determinant has a zero in the closed square of radius $10^{-12}$
about
\[
 (1.68430936268566911574800653125,\;
  5.40314724546132388337747818913).
\]
The corresponding original argument is approximately $17.13972580$.
Here is a separate certificate, independent of the large-index
argument. Put $L=\log w-i\pi/2$, $z=e^{aL}$, and
$D=E_{a,1}(z)^2-zE_{a,2}(z)E_{a,a}(z)$. For
$b\in\{1,2,a\}$ the term $z^j/\Gamma(aj+b)$ has first derivative
multipliers $jL-(j+b')\psi(aj+b)$ and $aj/w$ in $a$ and $w$,
respectively, where $b'=db/da$. The program
\texttt{schrodinger2.py} embedded in the source sums through $j=127$.
For $j\ge128$, all three series tails for each $b$ are bounded
using $3(j+1)^2 20^j/j!$ termwise. Indeed $|z|<20$, $|L|<4$,
$\Gamma(aj+b)\ge j!$ and $0<\psi(aj+b)<\log(2j+2)<j$.
The bounds have successive ratio below $1/2$ and their total is
less than $10^{-36}$.

For $P=(\Re D,\Im D)$ the fixed matrix
$B=10^{-3}\left(\begin{smallmatrix}896&-459\\10189&14195\end{smallmatrix}\right)$
is invertible. Arb enclosures at 176 bits verify that
$\|BP(a_0,w_0)\|_\infty<10^{-15}$ at the center and that
$\|I-BP'\|_\infty<1/100$ throughout the square.
Thus $x\mapsto x-BP(x)$ maps the square into itself and contracts
there. Its fixed point is a zero of the original determinant.


\section{A completely monotone squared Mittag--Leffler modulus}
\label{sec:schrod3}
Chorfi, Et-Tahri, Maniar and Yamamoto, \emph{Forward and backward problems for abstract time-fractional Schrödinger equations}, arXiv:2510.03600, Section~5, Conjecture~3, ask whether
\[
 f_\alpha(t)=E_{2\alpha,1}(-t)^2+tE_{2\alpha,1+\alpha}(-t)^2
\]
is completely monotone for $0<\alpha\leq 1/2$. Here $E_{\alpha,\beta}(z)=\sum_{k\ge0}z^k/\Gamma(\beta+\alpha k)$. We prove the stronger representation of $f_\alpha$ as the Laplace transform of a probability supported on $(0,\infty)$.

\paragraph{An auxiliary problem.}
Let $X=\sqrt{A}\,|G|$, where $G$ is centered Gaussian of variance $2$ and $A$ is any independent random variable supported on $(0,\infty)$. For an independent copy $Y$, is $X-Y$ a mixture of centered Gaussian laws? The answer is affirmative, including when the mixing law has atoms or has no moments.

\begin{proof}
For scale $a$ the positive-axis density is
\[
 h_a(x)=(\pi a)^{-1/2}\exp(-x^2/(4a)).
\]
Fix the scales $a,b$ of $X,Y$, respectively. The density of $X-Y$ at $x>0$ is
\[
 p_{a,b}(x)=\int_0^\infty h_a(x+y)h_b(y)\,dy
 =\frac{e^{-x^2/(4(a+b))}}{\sqrt{\pi(a+b)}}\,
 \operatorname{erfc}\!\left(\frac{x}{2\sqrt{a+b}}\sqrt{\frac ba}\right).
\]
This identity follows by completing the square. The conditional density at $-x$ is $p_{b,a}(x)$. Because the two scales are exchangeable, we may first symmetrize the conditional density by averaging these two expressions, then average in $a,b$.

For $z\ge0$,
\[
 \operatorname{erfc}(\sqrt z)=\frac2\pi\int_0^\infty
 \frac{e^{-z(1+s^2)}}{1+s^2}\,ds.
\]
Indeed, the two sides take value $1$ at zero and have derivative $-e^{-z}/\sqrt{\pi z}$ for $z>0$. Hence the symmetric density is a positive mixture of $e^{-\lambda x^2}$ with $0<\lambda<\infty$. This assertion holds first for each fixed $(a,b)$, and then for the averaged density by Tonelli's theorem. If $\nu$ is the resulting mixing measure, integration in $x$ gives $\int\sqrt{\pi/\lambda}\,\nu(d\lambda)=1$. Thus $\tau(d\lambda)=\sqrt{\pi/\lambda}\,\nu(d\lambda)$ is a probability measure. The Fourier transform of a normalized Gaussian yields
\[
 \mathbb E e^{i\xi(X-Y)}
  =\int_{(0,\infty)}e^{-\xi^2/(4\lambda)}\,\tau(d\lambda).
\]
This proves the assertion for every mixing law under consideration.
\end{proof}

\paragraph{Application to the original conjecture.}
For $0<\gamma<1$, take a positive stable random variable $S_\gamma$ with Laplace transform $e^{-u^\gamma}$. For completeness, existence follows from the Lévy--Khintchine representation
\[
 u^\gamma=\frac{\gamma}{\Gamma(1-\gamma)}
 \int_0^\infty(1-e^{-ux})x^{-1-\gamma}\,dx.
\]
Define $M_\gamma=S_\gamma^{-\gamma}$ and $M_1=1$. Integrating the Laplace transform against the gamma kernel gives, for every $p>0$,
\[
 \mathbb E M_\gamma^p=
 \frac{1}{\Gamma(\gamma p)}\int_0^\infty u^{\gamma p-1}e^{-u^\gamma}\,du
 =\frac{\Gamma(1+p)}{\Gamma(1+\gamma p)}.
\]
The final ratio also applies to $\gamma=1$. Monotone convergence in the exponential series gives finite moment generating function $E_{\gamma,1}(z)$ for all real $z\ge0$, hence the characteristic function is $E_{\gamma,1}(i\xi)$.

When $0<\alpha\le1/2$, take $A=M_{2\alpha}$ independent of $G$. The gamma duplication formula gives $\mathbb E|G|^k=k!/\Gamma(1+k/2)$ for integer $k\ge0$. Consequently $\sqrt A\,|G|$ and $M_\alpha$ have the same moments $k!/\Gamma(1+\alpha k)$, and the finite exponential moments identify their laws. Splitting the series into even and odd parts gives
\[
 E_{\alpha,1}(i\sqrt t)=E_{2\alpha,1}(-t)
     +i\sqrt t\,E_{2\alpha,1+\alpha}(-t).
\]
Taking squared moduli and applying the auxiliary result now shows that $f_\alpha(t)=\int e^{-t/(4\lambda)}\,\tau(d\lambda)$. In particular $(-1)^n f_\alpha^{(n)}(t)>0$ for every integer $n\ge0$ and every $t>0$. Cauchy--Schwarz in the integral also gives log-convexity.

\paragraph{Logarithmic interpolation for the evolution equation.}
In the Hilbert-space setting of equation (3.1) of the source, write $c_j=\langle u_0,\phi_j\rangle$ for the coefficients in the orthonormal eigenbasis, with $A\phi_j=-\lambda_j\phi_j$. Its solution formula (3.3) and Parseval's identity give
\[
 H(t):=\|u(t)\|^2=\sum_j |c_j|^2 f_\alpha(\lambda_j^2 t^{2\alpha}).
\]
Each term is a Laplace transform in $t$ because $e^{-d t^{2\alpha}}$ is the Laplace transform of a positive stable law rescaled by $d^{1/(2\alpha)}$ for $d>0$, with a point mass when $2\alpha=1$ or $d=0$. Summation therefore represents $H$ as a Laplace transform of a finite measure of mass $\sum_j |c_j|^2=\|u_0\|^2$. Hölder's inequality under this measure gives $H(t)\le H(0)^{1-t/T}H(T)^{t/T}$ for $0<t<T$. Taking square roots proves the authors' estimate (5.1), including its two endpoints by continuity,
\[
 \|u(t)\|\le\|u(0)\|^{1-t/T}\|u(T)\|^{t/T}
 \qquad(0\le t\le T,\quad 0<\alpha\le 1/2).
\]
The zero initial datum gives the inequality directly.


\section{Uniform positivity and first zeros in two fractional kernels}
\label{sec:cdh13}

Chaudhary, Diethelm and Hashemishahraki, \emph{On the separation of
solutions to fractional differential equations of order $\alpha\in(1,2)$},
Applied Numerical Mathematics \textbf{203} (2024), 84--96,
\href{https://doi.org/10.1016/j.apnum.2024.05.020}{doi:10.1016/j.apnum.2024.05.020},
formulate Conjectures~1 and~3 for the first positive zero $Z_\beta(\alpha)$
of $E_{\alpha,\beta}(-t^\alpha)$, with $\beta=\alpha$ and $1$,
respectively. The variable is $t$, as in their equation (17).
Writing $\beta=1+(1-\eta)(\alpha-1)$ with $\eta\in\{0,1\}$ makes
it possible to treat both problems in the same argument. Put
$S_0=2.93785$, $S_1=1.559$, $T_0=\pi$, and $T_1=\pi/2$.

\begin{theorem}\label{cdh13:result}
For $1\le\alpha\le2$ one has
$E_{\alpha,\beta}(-t^\alpha)>0$ on $0\le t\le S_\eta$.
For $\alpha=1+\delta\downarrow1$ let $L=\log(1/\delta)$ and now
allow every $0\le\eta\le1$, including parameters varying with $\delta$.
Uniformly throughout this interval,
\begin{equation}\label{cdh13:crossover}
 Z_\beta(1+\delta)
 =L+2\log L-\log(1+\eta L)+o(1).
\end{equation}
In particular the correction is $2\log L+o(1)$ when $\beta=\alpha$,
and $\log L+o(1)$ when $\beta=1$.
At the other endpoint, for $\eta\in\{0,1\}$,
\[
 \begin{gathered}
 Z_\beta(\alpha)=T_\eta+c_\eta(2-\alpha)+O((2-\alpha)^2),\\
 -0.810252<c_0<-0.810251,\qquad -0.129292<c_1<-0.129291.
 \end{gathered}
\]
\end{theorem}
The first assertion proves the exact positivity bounds in both part~(a)'s.
The leading coefficient of $\log(1/\delta)$ is exactly one in both
part~(b)'s. The printed numerical estimates $0.81$ and $0.275$ therefore
require correction. Because the source uses a relative error $1+o(1)$,
its additive fit parameter $d$ is unidentifiable at that precision.
The additive-error estimate here specifies the next diverging term.
For part~(c), the displayed intervals and the formula below specify the
slopes without assigning a tolerance to the authors' approximate values.

\paragraph{A moving-scale inverse Laplace estimate.}
Let $K_{\delta,\eta}(t)=t^{\beta-1}E_{\alpha,\beta}(-t^\alpha)$.
Its Laplace transform, obtained by integrating the defining series
first for real $p>1$, is
\[
 A_{\delta,\eta}(p)=\frac{p^{\eta\delta}}{1+p^{1+\delta}}.
\]
Use the principal logarithm. For $t\ge2$ and $0\le\delta\le 1/4$,
\begin{equation}\label{cdh13:uniform}
 K_{\delta,\eta}(t)=e^{-t}+\delta h_\eta(t)+O(\delta^2),\qquad
 \mathcal L h_\eta(p)=H_\eta(p)=
 \frac{[\eta+(\eta-1)p]\log p}{(1+p)^2}.
\end{equation}
Here the error is uniform even when $t$ tends to infinity.
To prove it, invert on $p=1/t+iy$. The Bromwich line can be shifted
there from real part $2$ because the transform is holomorphic in
the open right half-plane and the horizontal end segments vanish.
Taylor expand $A_{\delta,\eta}$ to first order in $\delta$ under
this inversion. Indeed, on this line
$|1+p^{1+v}|\ge c(1+|p|^{1+v})$ uniformly for $0\le v\le1/4$,
since $|(1+v)\arg p|\le5\pi/8$. Direct differentiation of its
quotient gives the following estimates for $B=\partial_v^2A_{v,\eta}$.
On $0<|y|\le1$ one has $|B|\le C(1+|\log|y||^2)$. On $|y|\ge1$
one has
\[
 |B|\le C\frac{1+\log^2(2+|y|)}{|y|},\qquad
 |\partial_y B|\le C\frac{1+\log^2(2+|y|)}{|y|^2}.
\]
For clarity, $B=A_{v,\eta}(\log p)^2[(\eta-U)^2-U(1-U)]$
where $U=p^{1+v}/(1+p^{1+v})$. Both $U$ and $1-U$ are bounded,
$\partial_pU=(1+v)U(1-U)/p$, and $A_{v,\eta}=O(1/|p|)$ at infinity.
These bounds justify the estimates and the integrated Taylor remainder.
On the two tails an integration by parts against $e^{iyt}$ gives
a uniform bound for the inverse integral of $B$. On the finite
interval the first estimate is integrable, and $e^{t\Re p}=e$.
This proves \eqref{cdh13:uniform}.

To determine $h_\eta(t)$, deform the inversion contour for $H_\eta$
to real part $-1/2$, splitting the contour at its cut on the negative
real axis. No pole at $-1$ is crossed. The jump of the logarithm gives
\[
 h_\eta(t)=-\int_0^{1/2}e^{-rt}
 \frac{\eta+(1-\eta)r}{(1-r)^2}\,dr+O(e^{-t/2}).
\]
The small indentation around zero has vanishing integral since
$H_\eta=O(1+|\log p|)$ there. The two vertical half-lines contribute
$O(e^{-t/2})$, by integration by parts away from their bounded
endpoints and the decay $H_\eta'=O((1+\log|p|)/|p|^2)$.
Expanding the integrand factor as $\eta+(1+\eta)r+O(r^2)$ yields
\[
 h_\eta(t)=-\frac{\eta}{t}-\frac{1+\eta}{t^2}+O(t^{-3})
  =-\frac{1+\eta t}{t^2}\bigl(1+O(t^{-1})\bigr).
\]
All constants are uniform for $0\le\eta\le1$, since
$\eta/(1+\eta t)\le1/t$.

Fix $v>0$ and define $t_\pm=L+2\log L-\log(1+\eta L)\pm v$.
Then $t_\pm/L=1+O(\log L/L)$ and
$(1+\eta t_\pm)/(1+\eta L)=1+O(\log L/L)$ uniformly. Therefore
\[
 \frac{e^{-t_\pm}t_\pm^2}{\delta(1+\eta t_\pm)}
       \longrightarrow e^{\mp v}
\]
uniformly. When the uniform remainder in \eqref{cdh13:uniform} is
divided by $\delta(1+\eta t)/t^2$, it is $O(\delta t^2)=o(1)$ for
$2\le t\le t_+$. The function $e^{-t}t^2/(1+\eta t)$ is decreasing
for $t>2$. Consequently $K$ is positive on $[M,t_-]$ and negative
at $t_+$ once $M$ is sufficiently large in terms of $v$ and $\delta$
sufficiently small. On $[0,M]$ the unscaled Mittag--Leffler series
tends uniformly to $e^{-t}$ as $\delta\to0$, by compact convergence
and uniform tail domination, uniformly also in $\eta$.
This confines the first zero to $[t_-,t_+]$ for all these parameters.
The arbitrariness of $v$ proves \eqref{cdh13:crossover}.
In the transition regime $\eta L\to\tau<\infty$ it supplies the
constant displacement $-\log(1+\tau)$.

\paragraph{The regular endpoint.}
Write $F(\alpha,t)=E_{\alpha,\beta}(-t^\alpha)$, $b=2-\eta$ and
$T=T_\eta$. At $\alpha=2$ the function is $\sin(t)/t$ for $\eta=0$
and $\cos t$ for $\eta=1$. Its first positive zero $T$ is simple.
The defining series and its first two derivatives in either variable
converge locally uniformly around $(2,T)$ by the gamma growth in
its terms. The implicit function theorem and compact positivity
before $T$ ensure that the local zero is the first zero. Hence
\[
 c_\eta=\frac{1}{F_t(2,T)}\sum_{k\ge0}
    \frac{(-1)^kT^{2k}}{\Gamma(2k+b)}
    \bigl[k\log T-(k+1-\eta)\psi(2k+b)\bigr].
\]
Here $F_t(2,T)=-1/\pi$ or $-1$ respectively. For $k\ge60$ the
summand has absolute value at most $6(k+1)^2 16^k/(2k)!$ using
$T<4$ and $0<\psi(s)<\log s$ for $s\ge2$.
The successive bounds have ratio below $1/2$. The first 60 terms
with this explicit remainder give the stated enclosing intervals.

\paragraph{Uniform positivity certificate.}
Let $X=t^\alpha$, $S=S_\eta$ and $m=1-\eta$. Define
$P(\alpha,X)=\sum_{k=0}^{60}(-X)^k/\Gamma(\alpha k+\beta)$.
Throughout $1\le\alpha\le2$ and $0\le X\le9$, the errors in the
value, the $X$ derivative, and the $\alpha$ derivative taken along
$X=S^\alpha$ are each strictly smaller than $10^{-13}$. Indeed,
for $k\ge61$ the corresponding three term magnitudes are at most
$4(k+1)^2 9^k/k!$. For the last bound use
$|k\log S-(k+m)\psi(\alpha k+\beta)|\le4(k+1)^2$ and monotonicity
of $\Gamma$ above $2$. The successive majorants have ratio below $1/2$.

We verify strict positivity at $X=S^\alpha$ and strict negativity
of the $X$ derivative on $[0,S^\alpha]$. The accompanying verifier
starts with $A=[1,2]$ and bisects it into dyadic intervals. On each
interval it computes a ball for the midpoint value at the right edge
and a ball for its derivative across $A$. The mean value theorem then
encloses the value along the edge. For the $X$ derivative, Horner
evaluation with interval inputs covers $[0,S^{\sup A}]$ after adaptive
bisection. Values of the reciprocal gamma function above $2$ and
of the digamma function on the positive axis are enclosed by monotonicity
and evaluation at both endpoints. Every comparison is a strict Arb
ball comparison after adding the tail bounds. At 112-bit precision,
the calculation succeeds on 567 final parameter intervals for $\eta=0$
(maximum depth 13) and on 37 for $\eta=1$ (maximum depth 8). This
proves positivity over the two full continua stated in the sources.

\paragraph{Earlier work and interpretation.}
For $\eta=1$ the leading logarithmic scale of the escaping zero
appears already in Gorenflo and Mainardi's earlier asymptotic balance
\cite[Section~3.2, Eq.~(3.36)]{GM}.
The present formula provides an additive $o(1)$ estimate including
the $\log L$ term. For $\eta=0$ the zero concerns the derivative of
the $\eta=1$ relaxation kernel, and its next coefficient is twice
as large. Formula~\eqref{cdh13:crossover} describes the transition
when $\beta$ moves between $1$ and $\alpha$ with the order. The certified bounds and slopes answer individually the
remaining questions in the two published conjectures. Historical
priority of the full collection is subject to further comparison.


\section{The drift of the zeros continuing from the sine function}
\label{sec:drift}
Let $Z_2(\alpha)$ be the least strictly positive $t$ for which
$E_{\alpha,2}(-t^\alpha)=0$. Conjecture~2 of Chaudhary, Diethelm and
Hashemishahraki \cite{CDH} includes the exact assertion
\[
 Z_2(2-\varepsilon)=\pi+\varepsilon+o(\varepsilon).
\]
This is part (b)(i) in the journal version and part (a)(i) in the arXiv version. We compute the drift at every
positive sine zero. The first already contradicts the stated coefficient.

\begin{theorem}\label{th:drift}
For each integer $k\ge1$, there is a unique real zero $t_k(\alpha)$
near $T=k\pi$ for $\alpha$ near $2$, and
\[
 t_k(2-\varepsilon)=T+
 \frac{T\operatorname{Ci}(T)+\operatorname{Si}(T)}{2}\,
 \varepsilon+O_k(\varepsilon^2).
\]
Here $\operatorname{Si}(x)=\int_0^x (\sin v)/v\,dv$ and
$\operatorname{Ci}(x)=\gamma+\log x+
 \int_0^x(\cos v-1)/v\,dv$. The first branch is $Z_2(\alpha)$.
Its coefficient exceeds $1.01$ and is approximately $1.041686$.
\end{theorem}
\begin{proof}
Put $Y(\alpha,t)=t E_{\alpha,2}(-t^\alpha)$. At $\alpha=2$ it equals
$\sin t$. The defining series and its first two derivatives in
$\alpha,t$ converge locally uniformly for $t>0$ and $\alpha$ near
$2$. This follows from Stirling's formula, since the derivatives add
only polynomial and logarithmic factors in the summation index on
each compact rectangle. The implicit function theorem at $T$ applies
because $\cos T\ne0$. For the first zero, choose a small fixed interval
about $\pi$ where the derivative in $t$ stays negative. On the
complementary interval from $0$ to its left endpoint,
$E_{2,2}(-t^2)=\sin(t)/t$ is strictly positive with limiting value one
at zero. Uniform convergence of the defining series there preserves
this positivity for $\alpha$ close to $2$.

Write $H_j=\sum_{l=1}^j1/l$ and
$A(t)=(t\cos t-\sin t)/2$. Termwise differentiation gives
\[
 \partial_\alpha Y(2,t)=
 \sum_{n\ge0}\frac{(-1)^n n t^{2n+1}}{(2n+1)!}
         (\log t+\gamma-H_{2n+1}).
\]
The same sum without the last factor equals $A(t)$.
Using $H_j=\int_0^1(1-u^j)/(1-u)\,du$, one obtains
\[
 \partial_\alpha Y(2,t)=(\log t+\gamma)A(t)
       -\int_0^1\frac{A(t)-A(tu)}{1-u}\,du.
\]
Termwise integration is justified by the convergent absolute series
weighted by $H_{2n+1}$. For $t=T$ substitute $v=T(1-u)$. The numerator
of the integral becomes
\[
 \frac{\cos T}{2}\{T(1-\cos v)+v\cos v-\sin v\}.
\]
Since $\int_0^T\cos v\,dv=0$, the preceding identities give
\[
 \partial_\alpha Y(2,T)=\frac{\cos T}{2}
       \bigl(T\operatorname{Ci}(T)+\operatorname{Si}(T)\bigr).
\]
Implicit differentiation with respect to $\varepsilon=2-\alpha$ gives
the displayed drift after division by $\cos T$.

For an elementary strict certificate set $p=157/50$ and $q=22/7$, so
$p<\pi<q$. Convexity of $1/x$ and the trapezoidal rule on $[5,m]$
give $\gamma\ge H_5-\log 5-1/10$ on letting $m\to\infty$.
The first four terms of the exponential series show
$e^{47/100}>5/p$, so $\gamma+\log\pi>H_5-1/10-47/100$.
For $j\ge0$ let $r_j=p$ if $j$ is even and $r_j=q$ if $j$ is odd.
The alternating series of the two integrals have decreasing term sizes
at $\pi$ (by the ratio test on consecutive term magnitudes).
Rational arithmetic therefore yields
\begin{align*}
 \operatorname{Ci}(\pi)&>
 H_5-\frac1{10}-\frac{47}{100}
 +\sum_{j=1}^{5}\frac{(-1)^j r_j^{2j}}{(2j)(2j)!}
 >\frac{60}{1000},\\
 \operatorname{Si}(\pi)&>
 \sum_{j=0}^3\frac{(-1)^j r_j^{2j+1}}{(2j+1)(2j+1)!}
 >\frac{1838}{1000}.
\end{align*}
Thus the first drift exceeds $(p\cdot0.060+1.838)/2>1.01$.
\end{proof}
The discrepancy concerns an exact coefficient in the published
conjunction. The following subsection establishes the positive assertions about the
threshold and the monotonicity of the first zero.
Implicit differentiation and differentiation of Mittag--Leffler series
are established tools. The result here identifies their consequence
for the stated first-zero law.


\subsection{A certified threshold, strict monotonicity and the square-root branch}
\label{cdh:fold-sec}

Write $F(a,t)=E_{a,2}(-t^a)$ and $Y_a(t)=tF(a,t)$. In the journal
version of Chaudhary--Diethelm--Hashemishahraki, Conjecture~2(a) asks
for positivity below a threshold $a_*$, and parts (b)(ii)--(iii) ask
for a square-root law above the threshold and strict decrease of the
first zero. Their cited phase diagram of Hanneken, Achar and Vaught
(2013, \href{https://doi.org/10.1155/2013/421685}{doi:10.1155/2013/421685},
Table~2) already lists the numerical threshold $1.59911520632302$.
We give rigorous enclosures, prove the global assertions, and determine
the square-root coefficient.

\begin{theorem}\label{cdh:fold-thm}
There are $a_*,t_*$ satisfying
\[
 |a_*-1.599115206323013|<10^{-13},\qquad
 |t_*-5.210687979825734|<10^{-11}
\]
such that $F(a,t)>0$ for every $1<a<a_*$ and $t\ge0$.
At $a_*$ there is exactly one positive zero $t_*$, of multiplicity two.
For $a_*<a\le2$ a first positive zero $Z_2(a)$ exists, and the
function $Z_2$ is strictly decreasing on $[a_*,2]$. As $a\downarrow a_*$,
\[
 Z_2(a)=t_* -\sqrt{\frac{-2F_a(a_*,t_*)}{F_{tt}(a_*,t_*)}}
              \sqrt{a-a_*}+O(a-a_*).
\]
The magnitude of this coefficient is between $4.787$ and $4.789$.
\end{theorem}

\paragraph{Certified compact estimates.}
Let $A$ and $T$ be the closed intervals in the first display. The
verifier establishes the following facts by outward-rounded arithmetic.
\begin{enumerate}
\item The rectangle $A\times T$ contains a simultaneous zero
$(a_*,t_*)$ of $(F,F_t)$, and the displayed bounds on the square-root
coefficient hold throughout that rectangle.
\item For every $a\in A$, one has $F>0$ when
$t\in[3,5.20]\cup[5.22,10]$, and $F_{tt}>0$ when $5.20\le t\le5.22$.
\item For every $a\in[1.5991,2]$, one has $F>0$ when $0\le t\le3$
and $F_a<0$ when $3\le t\le5.22$.
\end{enumerate}

The finite procedure uses $X=t^a$ and the series
$F=\sum_{k\ge0}(-X)^k/\Gamma(ak+2)$. Differentiate
term by term to obtain $F_t,F_{tt},F_a,F_{at}$. For $t\ge1$ their multipliers are
\[
 \frac{ak}{t},\quad \frac{ak(ak-1)}{t^2},\quad k(\log t-\psi(ak+2)),
 \quad \frac{k+ak^2(\log t-\psi(ak+2))}{t}.
\]
For $1.59\le a\le2$ and $1\le t\le10$, each term with $k\ge71$
including its multiplier is bounded in absolute value by
$20(k+2)^3 10^{1.59k}/\Gamma(1.59k+2)$. Indeed,
$10^{ak}/\Gamma(ak+2)$ decreases in $a$ on this range since
$\psi(1.59\cdot71+2)>\log10$ and $\psi$ is increasing.
One may use $\psi(ak+2)\le\log(ak+2)$ to bound the multipliers.
The successive majorants have ratio below $1/2$. This follows from
$\Gamma(x+s)/\Gamma(x)\ge x^s$ for $x>0,s\ge1$, by Jensen's inequality
applied to a gamma random variable. Twice the first majorant is
less than $10^{-36}$. The same bound holds for $F$ on $0\le t<1$
and for its $X$ derivative there, since for $k\ge71$ the maximum
of $k X^{k-1}/\Gamma(ak+2)$ occurs at the right edge of the full
$X$ interval, with $X\ge1$. This bound is added to every truncated
series result.

For the simultaneous zero, take the midpoint $x_0$ of $B=A\times T$,
write $P=(F,F_t)$ and let $M$ be the interval enclosure for $P'(B)$.
Take a fixed dyadic matrix $U$ from an inverse midpoint-Jacobian
calculation. The script checks $\det U\ne0$ and, in each coordinate,
\[
 |(UP(x_0))_i|+\sum_j|(I-UM)_{ij}|\operatorname{rad}(B_j)
                 <\operatorname{rad}(B_i).
\]
The continuous map $x\mapsto x-UP(x)$ therefore maps $B$ into its
interior. Brouwer's theorem yields the simultaneous zero there.
For the second item, the code checks $F$ using the mean value theorem
in $t$ and $F_{tt}$ by direct interval evaluation. It uses 49, 166
and 1 final time intervals on the three ranges, respectively.
For the third item, first check $F(a,3)>0$ and $F_X<0$ for
$0\le X\le3^a$ using 25 final parameter intervals and subdivisions
in $X$. For $F_a<0$ on the indicated rectangular strip it uses 129
final parameter intervals with subdivisions in $t$. The dyadic
subdivision and all precision and stopping rules are in the source.

\paragraph{Unbounded tail at the critical order.}
For $1<a<2$ Laplace inversion gives
\begin{align*}
 Y_a(t)={}&\frac2a e^{t\cos(\pi/a)}
          \cos\bigl(t\sin(\pi/a)-\pi/a\bigr)\\
  &+\frac{-\sin(\pi a)}\pi\int_0^\infty
     \frac{e^{-rt}r^{a-2}}{1+2r^a\cos(\pi a)+r^{2a}}\,dr.
\end{align*}
Indeed, its transform is $p^{a-2}/(1+p^a)$. Move the contour to the
negative real cut, using the principal power. At the two poles
$p=e^{\pm i\pi/a}$ the residues of the transform are $1/(ap)$.
The jump across the cut gives the displayed positive integral. The
large arc vanishes since the transform is $O(|p|^{-2})$, and the arc
around zero vanishes since $a>1$.

For $t\ge10$ keep only $0<r<1/t$ in that integral. It is bounded
below, including its coefficient, by $D_a t^{1-a}$ where
\[
 D_a=\frac{-\sin(\pi a)}{\pi(a-1)e(1+1/10)^2}.
\]
For all $a\in A$ the verifier checks $D_a10^{1-a}>(2/a)e^{10\cos(\pi/a)}$
and $\cos(\pi/a)+(a-1)/10<0$. Thus the exponential divided by the
power-law lower bound is decreasing on $[10,\infty)$ and already
below one at $10$. This proves strict positivity on the infinite tail.

\paragraph{Propagation between orders.}
The compact bounds and the tail imply $Y_{a_*}(t)\ge0$ everywhere.
Indeed, on $[5.20,5.22]$ the unique critical zero is a strict minimum
by strict convexity, and strict positivity holds off that interval.
Let $1<a<a_*$ and $q=a/a_*<1$. For the Laplace transforms $G_a$ of
$Y_a$ one has
\[
 G_a(p)=p^{-2(1-q)} G_{a_*}(p^q).
\]
Since $Y_{a_*}$ is nonnegative, $G_{a_*}$ is completely monotone.
Composition with the Bernstein function $p^q$ shows that $G_{a_*}(p^q)$
is the Laplace transform of a positive measure $\mu$. It satisfies
$G_{a_*}(p^q)\sim p^{-2q}$ as $p\to\infty$, so $\mu([0,u])>0$
for every $u>0$. Otherwise its transform would be $O(e^{-(p-1)u})$
by comparison with its value at $1$. The inverse transform of
$p^{-2(1-q)}$ is positive on all positive times. Its convolution
with $\mu$ is bounded below by a positive number locally about any
$t_0>0$: restrict the integral to $[0,t_0/2]$ and use times in
$[3t_0/4,5t_0/4]$. Uniqueness of Laplace transforms and continuity
of $Y_a$ give pointwise strict positivity.

At each $a>a_*$ the order-derivative bound gives $F(a,t_*)<0$.
Its first zero therefore belongs to $(3,t_*)$. If $a_*\le a_1<a_2\le2$,
evaluate $F(a_2,t)$ at the first zero for $a_1$. The strict
order-derivative bound again gives a negative value, proving
$Z_2(a_2)<Z_2(a_1)$. In particular the certified touching parameter
is the unique transition order.

As $a\downarrow a_*$ compactness and strict positivity before $t_*$
force $Z_2(a)\to t_*$. At the touching point, $F_t=0$, $F_{tt}>0$
and $F_a<0$. Put $a=a_*+s^2$ and $t=t_*+sv$. After division by $s^2$
the local equation extends analytically to $s=0$ with value
$F_a+F_{tt}v^2/2$. Both its roots in $v$ are simple. The implicit
function theorem and local strict convexity identify the branch with
the negative square root as the first zero and prove the displayed
expansion.


\section{The Guo--Qi conjecture for a normalized gamma-ratio derivative}
\label{sec:gqderiv}
B.-N. Guo and F. Qi, \emph{A class of completely monotonic functions
involving divided differences of the psi and tri-gamma functions and
some applications}, J. Korean Math. Soc. \textbf{48} (2011), 655--667,
\href{https://doi.org/10.4134/JKMS.2011.48.3.655}{doi:10.4134/JKMS.2011.48.3.655},
conjecture in Remark~4 the logarithmic complete monotonicity of their
function $\Phi_{s,t}$ from (36) for $|s-t|>1$ and of its reciprocal
for $|s-t|<1$. This includes the diagonal $s=t$. The 2009
preprint (arXiv:0903.1430) states the conjecture in Remark~3.

Write $x$ for the original argument translated by $\min(s,t)$,
so its entire domain is $x>0$, and put $r=|s-t|$. The original
function in these coordinates is $\Phi_r=P_r'$, where
\[
 P_r(x)=\begin{cases}
  \bigl[\Gamma(x+r)/\Gamma(x)\bigr]^{1/r}, & r>0,\\
  \exp\psi(x), &r=0 .
 \end{cases}
\]
In particular $\Phi_0(x)=\exp(\psi(x))\psi'(x)$.
We use the standard sign convention that a positive function
$\Phi$ is logarithmically completely monotone precisely when
$-(\log\Phi)'$ is completely monotone.

\begin{lemma}[Comparing a Laplace density with its canonical measure]
\label{gqderiv:compare}
Let $k$ be a strictly positive $C^2$ function on $[0,\infty)$ with
$k(0)=1$. Suppose $F(x)=\int_0^\infty e^{-xv}k(v)\,dv$ is finite
for every $x>0$, and $H=-F'/F$ is completely monotone.
Its positive representing measure has a continuous density $h\ge0$.
With $\ell=(\log k)'$, the exact identity
\[
 h(v)-k(v)=\int_0^v
  [\ell(v)-\ell(v-u)]k(v-u)h(u)\,du
\]
holds for $v\ge0$. Consequently $H-F$ is completely monotone for
log-convex $k$, and $F-H$ is completely monotone for log-concave $k$.
\end{lemma}
\begin{proof}
By the Bernstein--Widder theorem let $H$ be the Laplace transform
of the positive measure $\mu$. The product $FH=-F'$ and uniqueness
of Laplace transforms give
\[
 (k*\mu)(v):=\int_{[0,v]} k(v-u)\,d\mu(u)=v k(v)
\]
for almost every $v>0$. Both sides are right-continuous, so equality
holds for every $v\ge0$, with $\mu(\{0\})=0$.
Direct integration of $k(w)=1+\int_0^w k'(s)\,ds$ gives
\[
 (k*\mu)(v)=\mu([0,v])+\int_0^v(k'*\mu)(s)\,ds.
\]
All integrations on compact intervals are valid by local boundedness.
It follows that $\mu$ is absolutely continuous and its density
satisfies $h(v)=k(v)+vk'(v)-(k'*h)(v)$ almost everywhere.
The right hand side is continuous, so it gives a continuous
nonnegative version of $h$, with $h(0)=1$. Inserting
$v k'(v)=\ell(v)(k*h)(v)$ yields the identity and its two sign
consequences. Laplace transformation is legitimate since both $h$
and $k$ have finite transforms at every $x>0$.
\end{proof}

\begin{theorem}\label{gqderiv:theorem}
For every $r>1$, $\Phi_r$ is logarithmically completely monotone
on $(0,\infty)$. For every $0\le r<1$, $1/\Phi_r$ is logarithmically
completely monotone there. All the logarithmic derivative inequalities
are strict in these two ranges. At $r=1$, $\Phi_1=1$.
\end{theorem}
\begin{proof}
Put $F_r=(\log P_r)'$. For $r>0$ the trigamma series yields
\[
 F_r(x)=\frac{\psi(x+r)-\psi(x)}{r}
  =\frac{1}{r}\int_0^r\sum_{n\ge0}\frac{du}{(x+n+u)^2},
 \qquad
 F_0(x)=\sum_{n\ge0}\frac1{(x+n)^2}.
\]
These positive integral and sum representations put every $F_r$ in
the generalized Stieltjes class of order two. We invoke
G. K. Kristiansen, \emph{A proof of Steutel's conjecture},
Ann. Probab. \textbf{22} (1994), 442--452, in its standard consequence
that every nonzero generalized Stieltjes function of order two
is logarithmically completely monotone. For the explicit class formulation
see C. Berg, S. Koumandos and H. L. Pedersen, \emph{Nielsen's beta function
and some infinitely divisible distributions}, Theorem~2.1,
\href{https://arxiv.org/abs/1905.04131}{arXiv:1905.04131}.
In particular $H_r=-F'_r/F_r$ is completely monotone. This input
supplies positivity of the canonical measure. No positivity of the
difference $H_r-F_r$ follows without a further comparison.

Using $(x+a)^{-2}=\int_0^\infty v e^{-(x+a)v}\,dv$ gives $F_r$
as the Laplace transform of the normalized kernel
\[
 k_r(v)=\frac{1-e^{-rv}}{r(1-e^{-v})}\quad(r>0),
 \qquad k_0(v)=\frac{v}{1-e^{-v}},\quad k_r(0)=1 .
\]
For $v>0$ and $r>0$,
\[
 (\log k_r)''(v)=
 \frac{1}{4\sinh^2(v/2)}-\frac{r^2}{4\sinh^2(rv/2)} .
\]
The strict increase of $z\mapsto\sinh(z)/z$ on $(0,\infty)$ shows
that this is positive for $r>1$ and negative for $0<r<1$.
At $r=0$ the second term is $1/v^2$, giving the same negative
sign. Lemma~\ref{gqderiv:compare} therefore supplies complete monotonicity
of $H_r-F_r$ for $r>1$ and of $F_r-H_r$ for $r<1$.
Because $\Phi_r=P_rF_r$, these are exactly
$-(\log\Phi_r)'$ and $-(\log(1/\Phi_r))'$, respectively.
Strictness follows from $h(0)=1$, continuity, and the strict change
of $\ell$ on every positive interval. At $r=1$ the gamma recurrence
gives $P_1(x)=x$.
\end{proof}

\paragraph{Comparison with an earlier incompatible statement.}
Y.-C. Chen, T. Mansour and Q. Zou, \emph{On the complete monotonicity
of quotient of gamma functions}, Math. Inequal. Appl. \textbf{15}
(2012), 395--402, \href{https://doi.org/10.7153/mia-15-33}{doi:10.7153/mia-15-33},
Theorem~1, asserted failure of ordinary complete monotonicity of
$P_{s,t}$, identified in their (5) with the derivative $\Phi_{s,t}$.
Their claim for $|s-t|>1$ can be checked directly at $r=2$.
There $\Phi_2=(2x+1)/(2\sqrt{x(x+1)})$ and the inverse Laplace
density of $-(\log\Phi_2)'$ is
$(1+e^{-v})/2-e^{-v/2}=(1-e^{-v/2})^2/2\ge0$.
Hence $\Phi_2$ is logarithmically completely monotone and consequently
completely monotone. The preceding theorem establishes the entire
parameter range of the 2011 conjecture.


\section{The optimal reciprocal gamma-quotient threshold}
\label{sec:qwg}
Qi, Wei and Guo, \emph{Complete monotonicity of a function involving
the ratio of gamma functions and applications}, Banach J. Math. Anal.
\textbf{6} (2012), no.~1, 35--44,
\href{https://doi.org/10.15352/bjma/1337014663}{doi:10.15352/bjma/1337014663},
define, for $y>-1$, $b=y+1$ and $x>-b$,
\[
 h_{a,y}(x)=(x+b)^{-a}
       \exp\left(\frac{\log\Gamma(x+b)-\log\Gamma(b)}{x}\right).
\]
The quotient in the exponential takes the value $\psi(b)$ at $x=0$.
A positive function $h$ is logarithmically completely monotone (LCM)
when $(-1)^k(\log h)^{(k)}\ge0$ for every $k\ge1$.
The authors' Remark~4.5 conjectures that, if $b>1/2$, LCM of $1/h_{a,y}$
on its entire domain requires $a\le1/(2b)$. This statement also appeared
as Remark~5 in the 2009 Qi--Guo preprint, arXiv:0904.1101.

\begin{theorem}\label{qwg:main}
For every $b>0$, the reciprocal $1/h_{a,b-1}$ is LCM on $(-b,\infty)$
if and only if $a\le A(b)$. The sharp threshold equals $1$ when $b\le1/2$.
When $b>1/2$, define
\[
 B_b(s)=\frac{s}{(1+bs)(1-e^{-s})}.
\]
There is a unique positive root $t_b$ of
\[
       \int_{t_b}^{\infty} B_b'(s)\,\frac{e^{-bs}}s\,ds=0 .
\]
Then $A(b)=B_b(t_b)$, and
\[
       \frac{2}{1+2b}< A(b)<\min\{1,1/b\}.
\]
In particular, the conjectured bound fails for every $b>1/2$.
One explicit counterexample has $y=0$ and $a=2/3$.
\end{theorem}

\begin{proof}
Translate the full domain by putting $X=x+b>0$. Write
\[
 F_b(X)=\frac{\log\Gamma(X)-\log\Gamma(b)}{X-b}
        =\int_0^1\psi\bigl(b(1-u)+uX\bigr)\,du .
\]
This also gives the removable value at $X=b$.
The logarithm of the reciprocal is $a\log X-F_b(X)$.
Using the trigamma integral in $F_b'$ and then $t=us$ gives
\begin{align*}
 F_b'(X)
 &=\int_0^1 u\int_0^\infty
       \frac{s e^{-b(1-u)s-uXs}}{1-e^{-s}}\,ds\,du
   =\int_0^\infty e^{-Xt}K_b(t)\,dt,\\
 K_b(t)&=t e^{bt}\int_t^\infty\frac{e^{-bs}}{s(1-e^{-s})}\,ds .
\end{align*}
The exchange uses Tonelli's theorem for a nonnegative integrand.
Thus the negative logarithmic derivative of the reciprocal is the Laplace
transform of $K_b(t)-a$. Bernstein's theorem and uniqueness of Laplace
transforms imply that it is completely monotone exactly when
$K_b(t)\ge a$ everywhere. Indeed the signed density is continuous and
bounded as shown below; after any fixed positive exponential tilt both
it and a positive representing measure are finite. Consequently
$A(b)=\inf_{t>0}K_b(t)$.

Set $W_b(t)=e^{-bt}/t$. The last kernel is the normalized tail average
\[
 K_b(t)=\frac{1}{W_b(t)}
              \int_t^\infty B_b(s)\,[-W_b'(s)]\,ds.
\]
The bounded function $B_b$ has limits $1$ and $1/b$ at zero and
infinity. The same endpoint limits hold for $K_b$. At infinity this
follows from averaging on $[t,\infty)$. At zero the mass of
$[\varepsilon,\infty)$ in the normalized average is
$W_b(\varepsilon)/W_b(t)\to0$ for every fixed $\varepsilon>0$.

Logarithmic differentiation shows
\[
 \frac{B_b'(s)}{B_b(s)}
      =\frac1{s(1+bs)}-\frac1{e^s-1}.
\]
The power series of $(e^s-1-s)/s^2$ has positive coefficients,
starts at $1/2$, and increases strictly to infinity. Thus $B_b$ is
strictly increasing for $b\le1/2$, in which case its tail average
is also increasing with infimum $1$. For $b>1/2$, $B_b$ decreases
until a unique point $s_b$ and then increases. Moreover
$B_b'(0+)=1/2-b<0$. Integration by parts gives
\[
 K_b(t)-B_b(t)=\frac{1}{W_b(t)}
                \int_t^\infty B_b'(s)W_b(s)\,ds,
 \qquad
 K_b'(t)=\frac{-W_b'(t)}{W_b(t)}(K_b(t)-B_b(t)).
\]
The integral is positive for $t\ge s_b$. On $(0,s_b)$ it increases
strictly from $-\infty$ to a positive number. Its negative divergence
follows from $W_b(s)\sim1/s$ and $B_b'(0+)<0$. Hence exactly one
root $t_b$ exists. The kernel decreases strictly up to $t_b$ and
increases strictly thereafter, proving the characterization and the
upper bound for its minimum.

Finally $1/(1-e^{-s})>1$ and
$1/(1-e^{-s})=\tfrac12+\tfrac12\coth(s/2)>1/s+1/2$ for $s>0$.
Consequently
\[
 B_b(s)>\frac{\max\{s,1+s/2\}}{1+bs}.
\]
For $b>1/2$ the right side decreases on $(0,2)$ and increases on
$(2,\infty)$, so its minimum is $2/(1+2b)$. Applying this strict
bound at $t_b$ finishes the proof, since $2/(1+2b)>1/(2b)$.
\end{proof}

The auxiliary problem was to locate the minimum of a normalized
tail average with a one-valley integrand. Its derivative is determined
by a tail average of the derivative of that integrand. This observation
gives a unique critical point and a sharp implicit threshold. All
ingredients are elementary or classical. Historical priority for
the characterization is a separate question from the refutation.


\section{The two squared-tangent remainder questions}
\label{sec:tan-tails}

Write $a_j=[x^{2j}]\tan^2 x$ for $j\ge1$, and set
\[
  U_n(x)=\frac{\tan^2x-\sum_{j=1}^{n-1}a_jx^{2j}}{a_nx^{2n}},
  \qquad U_n(0)=1.
\]
This is $T_{2n-1}[\tan^2 x]$ in the notation of F.~Qi,
\emph{Theory of normalized remainders in Taylor series expansions},
\href{https://arxiv.org/abs/2512.15797}{arXiv:2512.15797v4},
Section~12.7.1. Its First Conjecture (X.-L. Liu and F. Qi,
\href{https://doi.org/10.1515/ms-2025-0051}{Math. Slovaca 75 (2025),
Remark~4.1}) asks whether $U_{n+1}/U_n$ is concave on
$(-\pi/2,\pi/2)$. The Second Conjecture (Y. Shuang, C.-Y. He and
F. Qi, \href{https://doi.org/10.7153/jmi-2025-19-80}{J. Math.
Inequal. 19 (2025), Section~4.1}) asks whether $U_n/U_{n+1}$ is
absolutely monotone on $(0,\pi/2)$ and completely monotone on
$(-\pi/2,0)$. They are separately posed questions about opposite
quotients. We give one proof for both.

\begin{theorem}\label{th:tanpaired}
For every $n\ge1$, the Maclaurin coefficients of $U_{n+1}/U_n$
at every strictly positive even degree are negative. Those of
$U_n/U_{n+1}$ at every such degree are positive. Both expansions
converge for $|x|<\pi/2$. Consequently both conjectures hold,
with strict concavity in the first and strict derivative signs in
the second.
\end{theorem}

\paragraph{A coefficient criterion.}
Suppose $D(y)=\sum_{j\ge0}d_j y^j$ is real analytic in a disk
of radius greater than $\rho>1$, with $d_0>0$, $D(\rho)=0$, and
$d_j\ge0$ for every $j\ge2$. Define
\[
 r=1/\rho,\qquad
 e=\frac{1}{d_0}\sum_{j\ge2}(j-1)d_j\rho^j.
\]
If $0\le e<1$ and $2r e/(1-e)<(1-r)^2$, then all coefficients
of degrees $k\ge2$ in $(1-y)^2/D(y)$ are positive and this series
is analytic on $|y|<\rho$.

Indeed, put $Q(y)=D(y)/(1-r y)$. Its constant is $d_0$, and
for $k\ge1$ its coefficient is
\[
 [y^k]Q(y)=-\sum_{j>k}d_j\rho^{j-k}\le0,
 \qquad Q(\rho)=-\rho D'(\rho)=d_0(1-e)>0.
\]
Thus the geometric expansion of $1/Q$ has nonnegative coefficients
$b_j$, and converges on $|y|\le\rho$, because the sum of the
absolute nonconstant coefficients of $Q/d_0$ weighted by $\rho^j$
equals $e<1$. In particular
\[
  \sum_{j\ge1}b_j\rho^j=b_0e/(1-e).
\]
If $v_j=[y^j](1/D)$, then $v_j=rv_{j-1}+b_j$ and
$v_j\ge b_0 r^j$. For $k\ge2$,
\[
 v_k-2v_{k-1}+v_{k-2}
   =(1-r)^2v_{k-2}-(2-r)b_{k-1}+b_k
   \ge b_0r^{k-2}\bigl[(1-r)^2-2r e/(1-e)\bigr]>0.
\]
This proves the criterion.

\begin{proof}[Proof of Theorem~\ref{th:tanpaired}]
Put $y=(2x/\pi)^2$ and $c=n+1/2$. The partial fractions for
$\sec^2 x$, or equivalently the derivative of the tangent
partial fractions, give
\[
 U_n(x)=F_n(y)/F_n(0),\qquad
 F_n(y)=\sum_{j\ge0}(c+j)y^j
             \sum_{m\ge1,\ m\ {\rm odd}}m^{-2n-2-2j}.
\]
The constant term of $\sec^2x$ does not enter since $n\ge1$.
For $n\ge2$, separate the $m=1$ pole and write
\[
 \begin{aligned}
 D(y)&=(1-y)^2F_n(y)=c-(c-1)y+(1-y)^2H(y),\\
 H(y)&=\sum_{m\ge3,\ m\ {\rm odd}}m^{-2n-2}
     \left(\frac{c}{1-uy}+\frac{uy}{(1-uy)^2}\right),\qquad u=m^{-2},
 \end{aligned}
\]
where $u$ inside the sum depends on $m$. The secondary term $H$ is
analytic on $|y|<9$. Its total weight $s=\sum_{m\ge3,\ m\ {\rm odd}}m^{-2n-2}$
satisfies
\[
 s\le3^{-2n-2}\left(1+\frac{3}{2(2n+1)}\right)
     < 2\cdot 9^{-(n+1)}
\]
by integral comparison with mesh two. As $u\le1/9$ and $c\ge5/2$,
\[
 \begin{split}
 H(0)&=cs,\\
 H(2)&\le(9c/7+18/49)s<3cs/2,\\
 H'(2)&\le(9(c+1)/49+36/343)s<cs/3.
 \end{split}
\]
The coefficients $d_j$ of $D$ are positive for $j\ge2$. In fact
each $u$ term contributes its weight times
$u^{j-2}(1-u)((c+j-2)(1-u)-2u)>0$.
Now $D(1)=1$ and $D(2)=2-c+H(2)<0$, so there exists
$1<\rho<2$ with $D(\rho)=0$.
The equation at this zero implies $c(\rho-1)\ge\rho$, whence
$1-1/\rho\ge1/c$. Moreover
\[
 \begin{split}
 d_0 e &= \rho D'(\rho)-D(\rho)+d_0\\
 &= H(0)+(\rho^2-1)H(\rho)+\rho(\rho-1)^2H'(\rho)
   <7cs.
 \end{split}
\]
Since $d_0\ge c$, we have $e<14\cdot9^{-(n+1)}<1/(2c^2+1)$.
For the last inequality check $n=2$ and note that
$2(c+1)^2+1<9(2c^2+1)$. Thus
$2r e/(1-e)<1/c^2\le(1-r)^2$. The criterion proves
$[y^k](1/F_n)>0$ for $k\ge2$.

For $n=1$, $1/U_1(x)=x^2\cot^2 x=x^2\csc^2 x-x^2$.
The usual partial fractions $\csc^2x=\sum_{j\in\mathbb Z}(x-j\pi)^{-2}$
give strictly positive coefficients in $x^2\csc^2x$ at every positive
even degree. Consequently the same strict positivity holds in $1/F_1$
from degree two in $y$ onward.

Finally $F_n(y)=F_n(0)+yF_{n+1}(y)$, so
\[
 R(y):=\frac{U_{n+1}(x)}{U_n(x)}
  =\frac{F_n(0)}{F_{n+1}(0)y}\left(1-\frac{F_n(0)}{F_n(y)}\right).
\]
It has constant 1 and all higher coefficients negative. Its expansion
converges in $|y|<1$ by the above representations. For $0\le y<1$
the ratio itself is positive. Therefore
$|1-R(y)|\le1-R(|y|)<1$ throughout this disk. The expansion
$1/R=\sum_{j\ge0}(1-R)^j$ then has strictly positive coefficients
of every degree. Substituting $y=4x^2/\pi^2$ and differentiating
proves the required strict curvature and derivative signs. Evenness
gives the alternating signs on the negative half-interval.
\end{proof}


\section{Sharp constants for He--Qi's beta-product majorization problem}
\label{sec:hqmajor}
C.-Y. He and F. Qi, \emph{Integral representations for specific Maclaurin
series and their applications to normalized remainder of exponential
function}, Open J. Math. Sci. \textbf{10} (2026), 296--305,
\href{https://doi.org/10.30538/oms2026.0287}{doi:10.30538/oms2026.0287},
pose Problem~3 for
\[
 S_{m,j}(x)=\sum_{h=0}^\infty
 \binom{h+m+j}{j}^{-1}\frac{x^h}{h!},
 \qquad m,j\in\mathbb N_0,\quad x\in\mathbb R .
\]
The beta integral in their Theorem~5 gives
$S_{m,j}(x)=j\int_0^1 t^m(1-t)^{j-1}e^{xt}\,dt$ for $j>0$,
while $S_{m,0}(x)=e^x$. We will also use the same integral to
extend the shape $j$ to positive reals.

\begin{theorem}\label{hq:sharp}
The two assertions of absolute monotonicity in He--Qi Problem~3,
for its functions $G$ and $H$ in (27)--(28), hold on the full real
line. For its remaining two functions in (29)--(30), write
$s=(p,q)$, $t=(a,b)$ in decreasing order, where $s$ strictly
majorizes $t$, so $p>a\ge b>q\ge0$ and $p+q=a+b$.
Put $P_s=S_{m,p}S_{m,q}$ and $P_t=S_{m,a}S_{m,b}$.
The largest constants for which $P_t-A P_s$ and $P_s-B P_t$,
respectively, are absolutely monotone on $\mathbb R$ are
\[
       A_* = \frac{a!\, b!}{p!\,q!},
       \qquad B_* = \frac{pq}{ab}.
\]
For $q=0$, the sharp value of $B$ is zero, whereas the source asks
for $0<B<1$. No positive $B$ has the stated property in this case.
For $q>0$ both optimal constants lie strictly between zero and one.
The formulas hold for every integer $m\ge0$.
\end{theorem}

\begin{proof}
Any transform $\int_{[0,2]} e^{xv}d\sigma(v)$ of a finite positive
measure is absolutely monotone on $\mathbb R$, by differentiation
under the integral. We first address
\[
 G_{m,n;k,l}=S_{m,n+l}S_{m+k,n}-S_{m+k,n+l}S_{m,n},
 \quad k,l\ge1 .
\]
For $n>0$, expand the two beta integrals and interchange the
two integration variables in half of the difference. The resulting
integrand, apart from a positive constant and the common factor
$e^{x(t+u)}(tu)^m[(1-t)(1-u)]^{n-1}$, is
\[
  (u^k-t^k)\{(1-t)^l-(1-u)^l\}\ge0 .
\]
For $n=0$ the difference equals $e^x(S_{m,l}-S_{m+k,l})$, whose
integrand contains the nonnegative factor $t^m(1-t^k)$.

The other assertion in (28) concerns
$H=S_{p,n+1}S_{q,n+1}-S_{a,n+1}S_{b,n+1}$.
After symmetrization its sign is given pointwise by
$t^pu^q+t^qu^p-t^au^b-t^bu^a\ge0$.
Indeed, for $t,u>0$ the two symmetrized sums are
$2(tu)^{(p+q)/2}\cosh(d\log(t/u)/2)$ for $d=p-q$ and $a-b$,
respectively. Since $p-q>a-b$, the inequality follows.

For the sharp constants, suppose first that $q>0$. In each product
integral substitute $v=1-t$, $w=1-u$ and group by $y=v+w$,
so $0<y<2$ and the exponential is $e^{x(2-y)}$. Put $z=v/y$.
For shapes $(i,j)$ the density on this slice is
\[
 ij\,y^{i+j-1}\int_0^1 W_y(z) z^{i-1}(1-z)^{j-1}\,dz ,
 \qquad
 W_y(z)={\bf1}_{\{yz\le1,\ y(1-z)\le1\}}
          [(1-yz)(1-y(1-z))]^m .
\]
We set $W_y=0$ outside the indicated set. On that set the base of
the power is $1-y+y^2z(1-z)$. Consequently $W_y$ is symmetric
and nonincreasing with $|z-1/2|$, also when $m=0$.

Use the symmetric probability measure proportional to
$z^{a-1}(1-z)^{b-1}+z^{b-1}(1-z)^{a-1}$ on $(0,1)$.
The ratio of the corresponding symmetrized $s$ and $t$ integrands is
\[
 R(z)=\frac{\cosh((p-q)c)}{\cosh((a-b)c)},\qquad
 c=\tfrac12|\log(z/(1-z))|.
\]
It is at least one and increases with $|z-1/2|$, since
$d\tanh(dc)$ is nondecreasing in $d\ge0$. Both the unweighted
integral of $R$ and the weighted one are finite because $q>0$.
Opposite monotonicity gives
$\mathbb E(W_y R)\le \mathbb E(W_y)\mathbb E(R)$, by integrating
$(W_y(z)-W_y(z'))(R(z)-R(z'))\le0$ against two independent copies.
Here $\mathbb E(W_y)>0$. Therefore the ratio of the $P_s$ to $P_t$
slice densities lies between
\[
 \frac{pq}{ab}
 \quad\hbox{and}\quad
 \frac{pq}{ab}\frac{\mathrm B(p,q)}{\mathrm B(a,b)}
 =\frac{\Gamma(p+1)\Gamma(q+1)}{\Gamma(a+1)\Gamma(b+1)} .
\]
Both claimed differences have nonnegative densities on every slice.
This proves their absolute monotonicity, with the displayed constants,
also for positive real shapes using gamma factors.

As $j\downarrow0$ the measures $jv^{j-1}\,dv$ on $[0,1]$
converge weakly to the atom at zero. Hence the continued $S_{m,j}$
and every one of its derivatives converge to $e^x$ for fixed $x$.
In the preceding argument take the spread shapes $(p-\varepsilon,
\varepsilon)$ and let $\varepsilon\downarrow0$, keeping their sum
fixed. This gives the asserted nonnegative differences at $q=0$,
with $0!=1$ and $B_*=0$.

It remains to prove optimality. For each $j>0$ the beta integral gives
\[
 S_{m,j}(x)\sim e^x\Gamma(j+1)x^{-j}\quad (x\to+\infty),
 \qquad S_{m,j}(-X)\sim j\,\Gamma(m+1)X^{-m-1}
       \quad(X\to+\infty).
\]
For the first estimate substitute $x(1-t)$ as variable and use
dominated convergence. For the second substitute $Xt$ on $0<t<1/2$,
where $(1-t)^{j-1}$ is bounded, and bound the remaining beta integral
by an exponentially small tail. The first estimate continues to hold
for $j=0$ by the exact formula $S_{m,0}=e^x$.
Thus $P_t/P_s\to A_*$ as $x\to+\infty$ and
$P_s/P_t\to B_*$ as $x\to-\infty$, the second ratio tending to zero
when $q=0$. Larger constants violate positivity of the undifferentiated
function. For $q>0$, strict concavity of $\log$ gives $pq<ab$,
and strict convexity of $\log\Gamma(j+1)$ gives
$\Gamma(a+1)\Gamma(b+1)<\Gamma(p+1)\Gamma(q+1)$.
This proves all statements of the theorem.
\end{proof}

\begin{corollary}[Finite products with real shapes]\label{hq:finite}
For real $m\ge0$, let $s,t$ be positive real vectors of the same finite
length, with $s$ majorizing $t$ in the usual weak-partial-sum sense
and with equal totals. For $P_v=\prod_i S_{m,v_i}$ the largest
constants such that $P_t-A P_s$ and $P_s-B P_t$ are absolutely
monotone on $\mathbb R$ are
\[
 A=\prod_i\frac{\Gamma(t_i+1)}{\Gamma(s_i+1)},
 \qquad B=\prod_i\frac{s_i}{t_i}.
\]
\end{corollary}
\begin{proof}
The Hardy--Littlewood--P\'olya two-coordinate characterization of
majorization (see A. W. Marshall, I. Olkin and B. C. Arnold,
\emph{Inequalities: Theory of Majorization and Its Applications},
2nd ed., Springer, 2011, Chapter~2, Sect.~B) gives a finite chain
from $s$ to $t$ by permutations and replacements
$(p,q)\mapsto(\lambda p+(1-\lambda)q,(1-\lambda)p+\lambda q)$,
where $0\le\lambda\le1$. Every intermediate shape is positive.
The slice proof applies to real $m\ge0$ and real positive shapes.
Multiply its two-coordinate inequalities in the order of absolute
monotonicity by the untouched factors, which are themselves
absolutely monotone by the beta integral. Compose the resulting
inequalities along the chain. This is valid because the cone of
absolutely monotone functions is closed under addition and
multiplication. The constants telescope to the displayed products.
The endpoint asymptotics in the proof above show sharpness for the
entire product, since the sums of all shapes agree.
\end{proof}

\paragraph{Qi's Fifth Conjecture.}
The ``Fifth Conjecture'' in F. Qi, arXiv:2512.15797,
Section~12.7.1.5, asks for absolute monotonicity of the three special
products (12.7.1)--(12.7.3). Its $G_n$ is $G_{0,n;1,1}$, its $H_n$
is the variance minor of the single beta weight with parameter $n$,
and its $J_n$ is the sharp $A_*$ inequality with $m=1$,
$(p,q)=(n+2,n)$, $(a,b)=(n+1,n+1)$. The above theorem therefore
settles the three clauses together, including the endpoint $n=0$
for $J_n$. Their original statements appear in Qi, Mathematics
\textbf{12} (2024), 2859, Remark~7,
\href{https://doi.org/10.3390/math12182859}{doi:10.3390/math12182859},
and in the cited paper of He and Qi. These specializations are
counted together with the general problem as one contribution.

For a direct view of their dependence, put
\[
 D_p(x)=\int_0^1\!\int_0^1 (t-u)^2[(1-t)(1-u)]^p e^{x(t+u)}\,dt\,du
 \qquad(p\ge0).
\]
With the chapter's indexing, its three left sides satisfy
\[
 G_n=\tfrac{n(n+1)}2 D_{n-1},\qquad
 H_n=\tfrac{n^2}2 D_{n-1}\quad(n\ge1),\qquad
 J_n=\tfrac{n+1}2D_n\quad(n\ge0).
\]
The first two follow by symmetrizing the beta integrals. For the
third, take $n>0$ and put $v=1-t$, $w=1-u$. After symmetrization
on each fixed $y=v+w$ the integrand, besides the factor
$(n+1)e^{x(2-y)}$, is
\[
 (1-v)(1-w)(vw)^{n-1}
       \bigl[vw-\tfrac n2(v-w)^2\bigr]
 =\tfrac12(1-v)(1-w)\frac{d}{dv}\bigl[(v-w)(vw)^n\bigr].
\]
Here $w=y-v$. At a slice endpoint one of $v,w,1-v,1-w$ vanishes.
Integration by parts therefore gives
$\frac12\int(v-w)^2(vw)^n\,dv$ on that slice.
The limit $n\downarrow0$ follows from the same weak convergence of
beta measures used above. These identities also prove strict
positivity of every derivative of the three special differences.

\paragraph{Earlier determinant formulas.}
Baricz and Ismail, \emph{Tur\'an type inequalities for Tricomi
confluent hypergeometric functions}, Remark~9,
\href{https://arxiv.org/abs/1110.4699}{arXiv:1110.4699}, display
squared-Vandermonde integral formulas for Kummer determinants, including
the kernel $D_p$ above, and credit earlier results of Ismail and
Laforgia. Although the cited monotonicity statement is phrased for
$x>0$, their displayed positive integral extends directly to real $x$.
The determinant positivity itself is therefore classical. The main
comparison here is between two majorized shape vectors on an entire
constant-sum slice, with sharp bounds in both directions and products
of arbitrary finite length.


\section{A large-power obstruction for normalized arcsine remainders}
\label{qi6:sec}

Feng Qi's \emph{Theory of normalized remainders in Taylor series
expansions}, arXiv:2512.15797, Section~12.7.1.6 (Sixth Conjecture),
asks whether $T_{2n+1}[(\arcsin x/x)^q]$ is logarithmically absolutely
monotone on $[0,1)$ for every positive integer $q$ and every integer
$n\ge0$. The author attributes the question to Section~1 of
\emph{Series and connections among central factorial numbers, Stirling
numbers, inverse of Vandermonde matrix, and normalized remainders of
Maclaurin series expansions}, Mathematics \textbf{13} (2025), 223,
\href{https://doi.org/10.3390/math13020223}{doi:10.3390/math13020223}.
The required signs are nonnegativity of all positive-order derivatives
of the logarithm. We give an exact counterexample and explain how
the truncation index and the exponent play different roles.

Put $y=x^2$, $B(y)=\arcsin(\sqrt y)/\sqrt y$ with $B(0)=1$, and
write $B(y)^q=\sum_{k\ge0}u_k(q)y^k$. For $s=n+1$ the original
normalized remainder is exactly $R_{q,s}(x^2)$, where
\[
 R_{q,s}(y)=\sum_{j\ge0}\frac{u_{s+j}(q)}{u_s(q)}y^j.
\]
Every denominator here is positive since $B$ has positive Taylor
coefficients. If $[y^j]\log R_{q,s}(y)<0$, then the derivative of
order $2j$ of its logarithm in $x$ is negative at zero and also
for all sufficiently small positive $x$.

\begin{theorem}\label{qi6:main}
The Sixth Conjecture is false. At $n=0$ and $q=56$ one has
\[
 \left.\frac{d^8}{dx^8}\log T_1[(\arcsin x/x)^{56}]\right|_{x=0}
 =-\frac{50349478}{86625}.
\]
For every fixed $n\ge0$ the conjectured property fails for all
sufficiently large integers $q$. Conversely, at each fixed positive
integer $q$ it holds for every sufficiently large $n$.
\end{theorem}

\paragraph{The finite counterexample.}
For $s=1$, $R_{q,1}(y)=(B(y)^q-1)/(qy/6)$. If
$a_k=\binom{2k}{k}/[4^k(2k+1)]$ are the coefficients of $B$, the
identity $B(B^q)'=qB'B^q$ gives $u_0=1$ and
\[
 ku_k=\sum_{j=1}^k ((q+1)j-k)a_j u_{k-j}.
\]
Put $r_k=u_{k+1}/u_1$ and $l_k=[y^k]\log R_{q,1}(y)$.
Logarithmic differentiation gives
$kl_k=kr_k-\sum_{j=1}^{k-1}jl_jr_{k-j}$. Substituting $k\le4$
from these two finite recurrences gives
\[
 l_4=-\frac{q^4}{3732480}+\frac{10181q^2}{16329600}
       +\frac{2497q}{226800}+\frac{765187349}{15717240000}.
\]
At $q=56$ this equals $-25174739/1746360000$, which proves the
stated derivative. The polynomial has just one positive root by
Descartes' rule of signs, and it is positive at $55$. Thus $56$ is
the first integer exponent detected by this particular coefficient.

\paragraph{Transfer from a bounded moment distribution.}
The following argument works for any real analytic base $B(y)$ at
zero with $B(0)=1$ and $B'(0)=a>0$. For fixed $k$, binomial expansion
of $(1+(B-1))^q$ gives $u_k(q)=(aq)^k/k!+O(q^{k-1})$. The term of
degree $k$ in $q$ comes from $(B-1)^k$. For each fixed $s\ge1$ the
coefficientwise limit as $q\to\infty$ is therefore
\[
 R_{q,s}(t/(aq))\longrightarrow
 M_s(t)=\sum_{j\ge0}\frac{s!}{(s+j)!}t^j=\mathbb E e^{tU},
\]
where $U$ has density $s(1-u)^{s-1}$ on $(0,1)$. The beta integral
verifies the last equality. The logarithms converge coefficientwise
too, since each coefficient depends polynomially on finitely many
coefficients of a series with constant term one.

For every $s$, at least one coefficient of $\log M_s$ is negative.
Indeed write $\log M_s(t)=\sum_{j\ge1}\kappa_jt^j/j!$. Here
$\kappa_2=v=s/[(s+1)^2(s+2)]>0$. If $\kappa_1,\dots,\kappa_{2N}$
were all nonnegative, exponentiation would give
\[
 \mathbb E U^{2N}\ge \frac{(2N)!}{2^N N!}\,v^N
                 \ge N!v^N \ge (Nv/e)^N.
\]
Taking $N=\lceil3/v\rceil$ contradicts $U\le1$. A strictly negative
coefficient consequently persists in the scaled tails at every
sufficiently large $q$. Positive scaling preserves its sign. For
$s=1$ this mechanism has a particularly short witness:
$M_1(t)=(e^t-1)/t$ and $[t^4]\log M_1=-1/2880$.

\paragraph{The opposite limit with the exponent fixed.}
The elementary identity
\[
 B(y)=\frac{\pi/2-\arcsin\sqrt{1-y}}{\sqrt y}
\]
holds near $y=1$ in its slit domain. The representation
$B(y)=\int_0^1(1-yu^2)^{-1/2}\,du$ gives analytic continuation on
$\mathbb C\setminus[1,\infty)$. For each fixed $q$ the convergent
local expansion in powers of $\sqrt{1-y}$ has first nonconstant
half-integer coefficient $-q(\pi/2)^{q-1}$ at exponent $1/2$.
Standard coefficient transfer for a slit-domain expansion (or
singularity analysis term by term through exponent $7/2$) gives
\[
 u_k(q)=\frac{q(\pi/2)^{q-1}}{2\sqrt\pi}\,k^{-3/2}
       \left(1+\frac{b_q}{k}+\frac{c_q}{k^2}+O_q(k^{-3})\right).
\]
Consequently the centered second difference of $\log u_k(q)$ is
$3/(2k^2)+O_q(k^{-3})>0$ for large $k$.

We recall the short reciprocal-series argument applied to any
sufficiently late tail. Its coefficients $v_j=u_{s+j}/u_s$ are
positive, $v_0=1$, and $v_j/v_{j-1}$ is nondecreasing. Set
$1-1/R_{q,s}(y)=\sum_{j\ge1}p_j y^j$. Then $p_1=v_1>0$ and
$p_j=v_j-\sum_{i=1}^{j-1}p_i v_{j-i}$. If $p_i\ge0$ for $i<j$,
bound $v_{j-i}\le(v_j/v_{j-1})v_{j-1-i}$ in this sum and use the
same convolution at $j-1$ to infer $p_j\ge0$. This is the classical
Kaluza sign argument. It also gives $0\le p_j\le v_j$. Therefore
for $0\le y<1$ the series $P(y)=\sum p_jy^j=1-1/R_{q,s}(y)$ takes
values in $[0,1)$. Expanding $\log R_{q,s}=\sum_{k\ge1}P^k/k$
now proves nonnegativity of all Taylor coefficients, with convergence
throughout $|y|<1$. Substituting $y=x^2$ proves logarithmic absolute
monotonicity through the entire original interval.


\section{Log-concavity of the normalized arctangent remainder}
\label{sec:qi-seventh}

The Seventh Conjecture in Section 12.7.1.7 of Feng Qi,
\emph{Theory of normalized remainders in Taylor series expansions}
(arXiv:2512.15797), following X.-L. Liu and F. Qi,
\emph{Logarithmically absolute monotonicity and series expansions of
normalized remainders for inverse sine and tangent functions}, asks
whether $T_{2n}[\arctan x]$ is logarithmically concave on $(0,1)$
for every $n\ge0$. Finite geometric division of $(1+t^2)^{-1}$ and
integration from zero to $x$ give
\[
 T_{2n}[\arctan x]=(2n+1)\int_0^1 \frac{t^{2n}}{1+x^2t^2}\,dt .
\]

\begin{theorem}\label{thm:qi-seventh}
For every real $a>0$, the function
\[
 F_a(x)=a\int_0^1\frac{t^{a-1}}{1+x^2t^2}\,dt
\]
satisfies $(\log F_a)''(x)<0$ for $0<x\le1$.
In particular the Seventh Conjecture holds for all its stated indices.
\end{theorem}
\begin{proof}
Put $d=2x^2/(1+x^2)$ and $r=a/[(1+x^2)F_a(x)]$.
Differentiating $x^a F_a(x)=a\int_0^x u^{a-1}/(1+u^2)\,du$
gives $xF_a'/F_a=r-a$ and then $xr'=r(a-d-r)$.
Consequently
\[
 -x^2(\log F_a)''=r^2-(a-1-d)r-a .
\]
We prove positivity through an inequality between a sum of two
nonnegative variables and their maximum.

Write $b=-\log x$, $z=\tanh b$, and $A=1+z$. Then $b,z\ge0$ and
$d=1-z$. For $s\ge0$ define
\[
 h(s)=\frac{\tanh(b+s)-z}{A},\qquad
 W=\int_0^\infty e^{-as}h'(s)\,ds,\qquad
 J=\int_0^\infty e^{-as}h(s)h'(s)\,ds .
\]
In the defining integral of $F_a$, use $t=e^{-s}$ and
$1/(1+e^{-2(b+s)})=(1+\tanh(b+s))/2$.
An integration by parts and the identity $1+x^2=2/A$ now yield
$(1+x^2)F_a=1+W$. Also $h'(0)=d$ and
$-h''=2(z+Ah)h'$. Integrating the latter identity against
$e^{-as}$ gives
\[
 d-aW=2zW+2AJ .
\]
The derivatives of $h$ are bounded and decay exponentially at
infinity, so the integrations are valid for every $a>0$.
Since $h(0)=0$ and $h'(s)>0$,
\begin{align*}
 W^2
 &=\int_0^\infty\!\int_0^\infty
    e^{-a(u+v)}h'(u)h'(v)\,du\,dv\\
 &<\int_0^\infty\!\int_0^\infty
    e^{-a\max(u,v)}h'(u)h'(v)\,du\,dv
 =2J .
\end{align*}
For the equality on the right, split the quadrant along its diagonal.
As $z\ge0$ and $A\ge1$ this implies
\[
 d-(a+z)W-W^2=zW+2AJ-W^2>0 .
\]
Finally $r=a/(1+W)$, and direct substitution gives
\[
 \frac{(1+W)^2}{a}\bigl(r^2-(a-1-d)r-a\bigr)
   =d-(a+z)W-W^2 .
\]
This proves the strict inequality, including at $x=1$.
\end{proof}

\begin{corollary}
The point $1$ is the largest possible common right endpoint for all
the integer-indexed remainders. At every fixed $x>1$ the logarithmic
second derivative of $T_{2n}[\arctan x]$ is positive for all sufficiently
large $n$.
\end{corollary}
\begin{proof}
The probability measures $a t^{a-1}dt$ on $[0,1]$ converge weakly to
the point mass at $1$ as $a\to\infty$. For fixed $x>0$ the kernel
$(1+x^2t^2)^{-1}$ and its first two derivatives in $x$ are continuous
on $0\le t\le1$. Therefore $(\log F_a)''(x)$ converges to
$2(x^2-1)/(1+x^2)^2$. Taking $a=2n+1$ gives the assertion.
\end{proof}

\paragraph{Source scope.}
Liu--Qi's monotonicity result for the same real remainder gives a sign
for the first derivative. Their logarithmic convexity result for the
imaginary-argument continuation concerns the normalized hyperbolic
arctangent remainder. The present assertion concerns the second
logarithmic derivative of the real arctangent remainder.


\section{Derivative signs of exponential remainders at every real center}
\label{sec:normalized-remainder-oscillation}
Y. Shuang, C.-Y. He, and F. Qi, J. Math. Inequal. \textbf{19} (2025),
1239--1251, \href{https://doi.org/10.7153/jmi-2025-19-80}{doi:10.7153/jmi-2025-19-80},
combine two assertions in Guess~3. Write $M_n(z)=T_n[e^z]$ for their
normalized remainder. They ask whether $M_n$ is logarithmically absolutely
monotone on $\mathbb R$ and whether $M_n/M_{n+1}$ is absolutely
monotone there, for all integers $n\ge0$. Here absolute monotonicity
requires nonnegative derivatives of all orders, and logarithmic
absolute monotonicity imposes this condition on derivatives of
$\log M_n$ of orders at least one. These assertions subsequently
appear separately in F. Qi, \emph{Theory of Normalized Remainders in
Taylor Series Expansions}, arXiv:2512.15797, Sections~12.7.1.3--4.
We retain the original joint count.

\begin{theorem}\label{th:exp-remainders}
Fix any real $x$ and integer $n\ge0$. For each of
$g=\log M_n$ and $g=M_n/M_{n+1}$, both
$\{k\ge1:g^{(k)}(x)>0\}$ and $\{k\ge1:g^{(k)}(x)<0\}$
have positive lower natural density. In particular both assertions
of Guess~3 fail at every order parameter and every real center.
\end{theorem}
\begin{proof}
Taylor's integral remainder and a change of variables give
\[
 M_n(z)=\frac{(n+1)!}{z^{n+1}}
       \left(e^z-\sum_{j=0}^n\frac{z^j}{j!}\right)
       =(n+1)\int_0^1e^{zv}(1-v)^n\,dv.
\]
Thus $M_n(0)=1$, $M_n$ is entire of order at most one, and it is
positive on the real line. It has a complex zero. Otherwise
Hadamard factorization would give $M_n(z)=e^{az+b}$, whereas
$(\log M_n)''(0)$ is the strictly positive variance of the
nondegenerate probability measure $(n+1)(1-v)^n\,dv$.
All zeros are simple. Indeed, at a zero $\zeta\ne0$ the derivative
of the expression in parentheses is $\zeta^n/n!$, so
$M_n'(\zeta)=(n+1)/\zeta$.

Consider either of the real meromorphic functions
\[
 f=M_n'/M_n,\qquad
 f=M_n/M_{n+1}=z/(n+2)+1/M_{n+1}.
\]
Each has a nonempty discrete set of simple nonreal poles. For the center
$x$, write $R$ for the distance to the nearest poles and
$x+R e^{i\theta_j}$ for all poles at that distance.
Writing $c_j$ for their nonzero residues, choose $r>R$ which still
excludes every farther pole, including its boundary. Subtract these
principal parts and use Cauchy's coefficient estimate on the circle
of radius $r$. For $k\ge0$ this gives
\[
 \frac{R^{k+1}}{k!} f^{(k)}(x)
  =v_k+O\bigl((R/r)^k\bigr),\qquad
 v_k=-\sum_j c_j e^{-i(k+1)\theta_j}.
\]
The sequence $v_k$ is real by conjugation. Its arithmetic mean tends to zero
because no $\theta_j$ is zero modulo $2\pi$. Its mean square tends to
$V=\sum_j|c_j|^2>0$ by the geometric-sum identity for distinct
frequencies. Since $|v_k|\le L=\sum_j|c_j|$, the averages
of $|v_k|$ have lower limit at least $V/L$. The positive and negative
parts therefore both have mean lower limit at least $V/(2L)$.
Taking $\delta=V/(4L)$ shows that each of $v_k>\delta$ and
$v_k<-\delta$ has lower density at least $V/(4L^2)$. The error in
the coefficient estimate tends to zero. Both sign sets for the
derivatives of $f$ thus have positive lower density. For $f=M_n'/M_n$
these are precisely derivatives of $\log M_n$ with the index shifted
by one, which preserves lower density.
\end{proof}
The argument uses classical factorization and coefficient estimates.
The positive-density conclusion quantifies the failure of the original
conjecture for every center and order parameter, with
density bounds that may depend on those choices.
Nearest-singularity explanations of high-cumulant oscillations have
classical precedents, including C. Flindt et al., \emph{Universal
oscillations in counting statistics},
\href{https://doi.org/10.1073/pnas.0901002106}{doi:10.1073/pnas.0901002106}.
The application here simultaneously treats all the exponential
remainders, all their adjacent quotients, and every real center.


\section{Adjacent moment tails for inverse sine and inverse tangent}
\label{sec:moment-ratios}
The Eighth and Ninth Conjectures in Section~12.7.1 of Qi's chapter
ask about the ratio $T_{2n}[f]/T_{2n+2}[f]$ for $n\ge0$.
They ask for concavity on $(-1,1)$ for $f=\arcsin$, and for increase
on $(0,\infty)$ for $f=\arctan$. The following auxiliary problem
treats both signs of the argument for one measure. It also gives
an all-orders extension of the second assertion.

Let $\nu$ be a probability measure on $[0,1]$, with
$m=\int s\,d\nu>0$. Define its size-biased law
$d\nu_1(s)=s\,d\nu(s)/m$, let $m_1=\int s\,d\nu_1$, and put
\[
 F(z)=\int\frac{d\nu(s)}{1-zs},\qquad
 F_1(z)=\int\frac{d\nu_1(s)}{1-zs},\qquad R(z)=F(z)/F_1(z).
\]
\begin{proposition}\label{th:sizebias}
For some finite positive measure $\rho$ on $[0,1]$,
\[
 R(z)=1-(m_1-m)z-z^2\int\frac{d\rho(s)}{1-zs},\qquad z<1.
\]
If $\nu$ is nondegenerate, $R(x^2)$ is strictly concave on
$(-1,1)$. The function $w\mapsto R(-w)$ is then positive and strictly
increasing for $w\ge0$, with a completely monotone first derivative.
\end{proposition}
\begin{proof}
In $L^2(\nu_1)$ let $A$ be multiplication by $s$, let $e=1$,
let $Q$ project orthogonally onto $e^\perp$, and write
$C=QA|_{e^\perp}$, $v=QAe$. Since $0\le C\le I$, the
Schur complement of $I-zC$ in $I-zA$ gives
\[
 F_1(z)^{-1}=1-zm_1-z^2\langle v,(I-zC)^{-1}v\rangle .
\]
The last inner product is the transform of the spectral measure
of $v$ for $C$. The identity $F=1+mzF_1$ proves the representation.
This uses only the classical resolvent identity and spectral theorem.

Since $m_1-m=\operatorname{Var}_\nu(s)/m$, the coefficient of $z$
is strictly negative for nondegenerate $\nu$. Every coefficient
of higher degree is nonpositive, proving the strict concavity
after setting $z=x^2$.
For $w\ge0$, $R(-w)=m/\mathbb E_w[s]$, where
the probability in $\mathbb E_w$ has density proportional to
$(1+ws)^{-1}$ relative to $\nu$. Differentiation gives
\[
 \frac{d}{dw}\mathbb E_w[s]
 =-\operatorname{Cov}_w\!\left(s,\frac{s}{1+ws}\right)<0,
\]
because both functions of $s$ are strictly increasing. For $j\ge2$
the representation gives
\[
 (-1)^{j-1}\frac{d^j}{dw^j}R(-w)
 =j!\int\frac{s^{j-2}}{(1+ws)^{j+1}}\,d\rho(s)\ge0 .
\]
This proves the derivative signs on the entire positive half-line.
\end{proof}

For arcsine the coefficients of $\arcsin x/x$ in $x^2$ are
$a_j=4^{-j}\binom{2j}{j}/(2j+1)$. These are the moments of
$S=VU^2$, where $U$ is uniform on $[0,1]$ and $V$ is independent
with the arcsine law $dv/(\pi\sqrt{v(1-v)})$.
Use as $\nu$ the law of $S$ weighted by $S^n/a_n$.
Then $T_{2n}[\arcsin x]=F(x^2)$, and advancing $n$ by one
replaces $F$ by $F_1$. Proposition~\ref{th:sizebias} proves the
Eighth Conjecture. For arctangent use $S=U^2$, whose moments
are $(2j+1)^{-1}$. Then $T_{2n}[\arctan x]=F(-x^2)$.
The identity holds on the full real line by the geometric remainder
under the integral for $\arctan x/x$. It gives the strictly increasing
quotient in the Ninth Conjecture for every $x>0$.

\paragraph{Comparison with earlier results.}
For the Ninth Conjecture the hypergeometric representation is
$T_{2n}[\arctan x]={}_2F_1(1,n+1/2;n+3/2;-x^2)$.
Karp and Sitnik, \emph{Inequalities and monotonicity of ratios for
generalized hypergeometric function}, J. Approx. Theory 161 (2009),
337--352, \href{https://doi.org/10.1016/j.jat.2008.10.002}{Theorem~1},
already prove the required monotonicity on taking the upper and lower
parameter shifts equal to one and inverting the ratio. For the Eighth
Conjecture, the arcsine moment coefficients are strictly log-convex.
Kaluza's classical reciprocal-series criterion gives nonpositive
nonconstant coefficients of $1/F_1$. The identity $F/F_1=mz+1/F_1$
and $m_1>m$ then prove the claimed concavity directly. These two items
are excluded from the ten in Section~1.


\section{Maximum-degree neighborhoods and Graffiti.pc Conjecture 165}
\label{wow165:sec}

Let $G$ be a finite simple connected graph. Write $L(G)$ for the maximum
number of leaves in a spanning tree, $M$ for its set of maximum-degree
vertices, and $t(G)=|N_G(M)\setminus M|$ for the number of distinct
vertices outside $M$ adjacent to it. The 8 August 2005 entry for
Graffiti.pc Conjecture~165 asks whether $L(G)\ge\lceil\sqrt{2t(G)}\rceil$
\cite{wow165:registry}. The entry discusses a summer project of Craig
Foster which asks for a universal constant $c$ with $L(G)\ge c\,t(G)$.
His examples require $c\le1/4$.
\begin{theorem}\label{wow165:theorem}
Every finite simple connected graph of maximum degree at least three satisfies
\[
 t(G)\le4\bigl(L(G)-2\bigr).
\]
The coefficient $4$ is sharp. Conjecture~165 holds on its entire
original domain, and the largest constant in Foster's multiplicative
inequality is $1/4$.
\end{theorem}

\paragraph{A component problem.}
Retain exactly the edges incident with $M$ and discard isolated vertices.
In each resulting component the vertices outside $M$ are independent.
The quantity $L-2$ gives a budget for assembling solutions on these
components, since for every tree $T$ on at least two vertices,
\begin{equation}\label{wow165:budget}
 \#\mathrm{leaves}(T)-2=\sum_{u\in V(T)}(\deg_T u-2)_+,
 \qquad x_+=\max(x,0).
\end{equation}
In particular, extending vertex-disjoint trees to a spanning tree of a
connected graph preserves each selected vertex's contribution to the
right-hand side or increases it. We quantify how large the external
part of a single component can be under this budget.

\begin{lemma}\label{wow165:component}
Let $H$ be a finite simple connected graph with vertex partition
$A\cup B$, where $A$ is nonempty, $B$ is independent, all vertices of
$A$ have a common degree $D\ge3$, and the degrees in $B$ are less than
$D$. Then $|B|\le4(L(H)-2)$.
\end{lemma}
\begin{proof}
Set $h=|A|$, $t=|B|$, $q$ equal to the number of components of $H[A]$,
and $g=L(H)-2$. The full star about a vertex in $A$ extends to a spanning
tree without decreasing the contribution of its center in
\eqref{wow165:budget}, so $g\ge D-2\ge1$. If $H$ is a tree,
\eqref{wow165:budget} gives $g\ge h(D-2)$. Also $t\le Dh$ by counting
neighbors, and $D\le3(D-2)$, proving $t\le3g$ in this case.

Suppose $H$ contains a cycle. Extend a spanning forest of $H[A]$ to
a spanning tree $T$ of $H$. All $h+t-1$ edges of $T$ touch $A$, and
the $h-q$ edges in its prescribed forest touch $A$ twice. Thus
\[
 g\ \ge\sum_{u\in A}(\deg_T u-2)
   \ \ge\ (h+t-1)+(h-q)-2h=t-q-1.
\]
This gives $t\le g+q+1$. Every degree-two vertex belongs to $B$, so
there is no pair of adjacent degree-two vertices. Bankevich and Karpov's
Theorem~2 \cite{wow165:BK}, applied with girth lower bound $3$ and
maximum degree-two chain length $1$, gives
\[
   L(H)\ge (h+t-3)/6+2 .
\]
Here its coefficient is $(3-2)/((3-1)(1+2))=1/6$. Consequently
\begin{equation}\label{wow165:est}
  t\le g+1+q\le g+1+h,\qquad h+t\le6g+3,\qquad 2t\le7g+4 .
\end{equation}
For $g\ge4$ this implies $t\le4g$, while $g=3$ gives $t\le12$ by
integrality.

For the two remaining values we record a packing observation. Two
nonadjacent degree-$D$ vertices with disjoint neighborhoods give two
disjoint full stars. Extending their union to a spanning tree and using
\eqref{wow165:budget} yields at least $2D-2$ leaves. If the two centers
have $s\ge1$ common neighbors, join the centers by one of these
neighbors and attach each remaining vertex in the union of their
neighborhoods to one incident center. When both centers have other
assigned neighbors, this tree has $2D-s-1$ leaves. Such an assignment
exists since $D\ge3$: if the two neighborhoods are identical, give each
center a different remaining neighbor; otherwise each has a private
neighbor. Extending the tree gives $L(H)\ge2D-s-1$.

If $g=1$ and $t>4$, \eqref{wow165:est} forces $t=5$ and $q\ge3$.
We have $D=3$. Take one center in each of three distinct components of
$H[A]$. Every pair of these nonadjacent centers must have at least
two common neighbors by the packing observation and $L(H)=3$.
A common neighbor lies in $B$, since otherwise the centers belong to
the same component of $H[A]$. Each vertex of $B$ has degree at most
two and serves at most one pair. The three pairs thus require six
distinct vertices of $B$, contradicting $t=5$.

If $g=2$ and $t>8$, \eqref{wow165:est} gives $t=9$ and
$6\le q\le h\le6$. Thus $A$ is independent and $D\in\{3,4\}$.
When $D=4$, the packing observation and $L(H)=4$ require at least
three common neighbors for every pair of centers. There are fifteen
pairs. On the other hand,
\[
 \sum_{b\in B}\binom{\deg_H b}{2}\le\sum_{b\in B}\deg_H b=24,
\]
since $\deg_H b\le3$. This contradicts the required count of at least
$45$. For $D=3$ the degree sum over $B$ is $18$, and each of its nine
vertices therefore has degree exactly two. Suppress them to obtain a
connected loopless cubic multigraph on the six centers with nine
edges. Take a spanning tree of the multigraph and fix a leaf $a$ in
it. Subdivide each of its five edges as in $H$. For each of the four
unused edges, attach its subdivision vertex as a leaf to an endpoint
different from $a$. Such an endpoint exists because the multigraph is
loopless. The resulting spanning tree of $H$ has those four leaves
and $a$, contradicting $L(H)=4$. This completes the lemma.
\end{proof}

\begin{proof}[Proof of Theorem~\ref{wow165:theorem}]
Let the maximum degree of $G$ be $D\ge3$. Form the spanning subgraph
with all edges incident to $M$. Its nontrivial components $H$ satisfy
Lemma~\ref{wow165:component} with $A=V(H)\cap M$ and $B=V(H)\setminus M$.
Their $B$ sets partition $N_G(M)\setminus M$. In each $H$ choose a
spanning tree with $L(H)$ leaves. The trees form a forest in $G$,
and extend to a spanning tree of $G$. Equation~\eqref{wow165:budget}
gives
\[
 t(G)=\sum_H |V(H)\setminus M|
 \le4\sum_H(L(H)-2)\le4(L(G)-2).
\]
In particular $L(G)^2-2t(G)\ge(L(G)-4)^2\ge0$. For maximum degree at
most two the possibilities are a singleton, a path and a cycle.
The singleton, $K_2$, and every cycle have $t=0$. A path on at least
three vertices has $t=2$ and $L=2$. These cases also satisfy the
stated square-root inequality and $L\ge t/4$.

For sharpness, the original annotation credits Foster with the
parameters $L=k+2$, $t=4k$. An explicit realization described in
VMRC~[066] is as follows. For each cyclic index $i=1,\ldots,k$ take
two length-two paths from $a_i$ to $b_i$ and one length-three path
from $b_i$ to $a_{i+1}$, with all internal vertices distinct. The
$2k$ endpoints have degree three. The $4k$ internal vertices comprise
their external neighborhood. Every spanning tree breaks each of the
$k$ vertex-disjoint four-cycles, so at most one endpoint of each pair
retains degree three. Thus $L\le k+2$ by \eqref{wow165:budget}, and
the bound just proved gives equality. Consequently $L/t$ tends to
$1/4$, as required.
\end{proof}

\paragraph{Historical scope.}
The 2006 UHD conference abstract by Foster \cite{wow165:Foster}
reports the resolution of several Graffiti.pc conjectures and
discusses this specific pair of invariants. It does not identify those
conjecture numbers or state the inequalities and hypotheses. Its scope
therefore leaves a substantive priority question. Corpus [066] treats the case where $M$ is
independent, $D\in\{3,4,5\}$ and the remaining degrees are at most
two. The proof above gives the result for arbitrary connected graphs.
Historical priority for the full inequality or its method has not
been established.


\section{A sharp two-shift classification for the Qi--Luo question}
\label{sec:qiluo}
Qi's 2011 preprint (arXiv:1104.4442, Remark~5) and the revised article
by F.~Qi and Q.-M.~Luo, \emph{Complete monotonicity of a function
involving the gamma function and applications}, Period. Math.
Hungar. \textbf{69} (2014), 159--169,
\href{https://doi.org/10.1007/s10998-014-0056-x}{doi:10.1007/s10998-014-0056-x},
Remark~4.6, ask for the logarithmic complete monotonicity of their
two-parameter function (4.5). For a correction amplitude $c>0$ put
\[
 G^{(c)}_{\lambda,\mu}(x)=
 \frac{e^x\Gamma(x+1)}{(x+\mu)^{x+\mu}}
 \exp\!\left(\frac{c}{x+\lambda}\right),
 \qquad x>\max(0,-\lambda,-\mu).
\]
Both shifts are arbitrary real numbers. The source fixes $c=1/24$.
We ask in addition which correction amplitudes allow the region of
logarithmic complete monotonicity to be read from the two endpoint
rates of a Laplace kernel.

\begin{theorem}\label{th:qiluo}
For every $c\ge1/24$ and all real $\lambda,\mu$, the function
$G^{(c)}_{\lambda,\mu}$ is logarithmically completely monotone on its
full domain if and only if
\[
 \mu\ge \max\{1/2,\min(1,\lambda)\}.
\]
The number $1/24$ is the smallest positive amplitude for which this
criterion holds for every pair of shifts. Taking $c=1/24$ answers the
two-parameter question in Remark~4.6.
\end{theorem}
\begin{proof}
Logarithmic complete monotonicity is equivalent to complete
monotonicity of the negative logarithmic derivative. The latter is
\[
 \log(x+\mu)-\psi(x+1)+\frac{c}{(x+\lambda)^2}
    =\int_0^\infty e^{-xt}\frac{K(t)}{t}\,dt,
 \quad K(t)=\frac{t}{e^t-1}-e^{-\mu t}+ct^2e^{-\lambda t}.
\]
This follows by subtracting the integral for the digamma function
from Frullani's formula for the logarithm and adding the elementary
Laplace integral for the last term. The density is bounded near zero
and integrable after multiplying by $e^{-xt}$ for every admissible
$x$. Its continuity, Bernstein's representation theorem, and uniqueness
of the Laplace transform after a fixed admissible exponential tilt
give the criterion $K(t)\ge0$ for every $t>0$.

Since $K(t)=(\mu-1/2)t+O(t^2)$, necessity gives $\mu\ge1/2$.
If $\mu<\min(1,\lambda)$, multiplying $K(t)$ by $e^{\mu t}$ gives
limit $-1$ at infinity. These are the two necessary endpoint conditions.

For sufficiency it is enough to treat $c=1/24$ and the smallest
allowed $\mu$, since increasing either only increases $K$. For $u>0$,
\[
 \frac{u}{\sinh u}+\frac{u^2}{6}\ge1.
\]
Indeed the series of $(1-u^2/6)\sinh u$ has coefficient one at
degree one and coefficient
$[1-(2j)(2j+1)/6]/(2j+1)!\le0$ at degree $2j+1$ for $j\ge1$.
Put $u=t/2$, so $e^{t/2}t/(e^t-1)=u/\sinh u$.
For $\lambda\le1/2$ take $\mu=1/2$, and observe
\[
 e^{\mu t}K(t)=\frac{u}{\sinh u}-1
       +e^{(1/2-\lambda)t}\frac{u^2}{6}\ge0.
\]
For $1/2\le\lambda\le1$ take $\mu=\lambda$, giving
\[
 e^{\mu t}K(t)=e^{(\lambda-1/2)t}\frac{u}{\sinh u}
                    -1+\frac{u^2}{6}\ge0.
\]
For $\lambda\ge1$ take $\mu=1$. Here
$e^t t/(e^t-1)=t/(1-e^{-t})>1$.
Finally, at $\lambda=\mu=1/2$ the normalized kernel is
$u/\sinh u-1+ct^2=(c-1/24)t^2+O(t^4)$,
so every smaller amplitude fails for this pair.
\end{proof}

\paragraph{The reciprocal.}
Although Remark~4.6 asks explicitly about $G$ itself, its reciprocal
at the source amplitude has the following sharp variational threshold.
It is logarithmically completely monotone if and only if $\mu\le1/2$
and
\[
 \lambda\ge\Lambda(\mu):=
 \max_{t>0}\frac1t
     \log\frac{t^2}{24\bigl(e^{-\mu t}-t/(e^t-1)\bigr)}.
\]
Indeed the requisite kernel inequality is exactly $K(t)\le0$.
The expansion at zero requires $\mu\le1/2$. Under that condition
the parenthesized difference is positive because $u/\sinh u<1$.
Taking logarithms gives the formula with a supremum. The expression
in the maximum tends to $-\infty$ at zero if $\mu<1/2$, and to $1/2$
there if $\mu=1/2$. At infinity it is
$\mu+(2\log t-\log 24+o(1))/t$, so it tends to $\mu$ from above.
These endpoint facts imply the maximum is attained at a finite
positive point. Qi's one-shift threshold at $\mu=1/2$ is a predecessor
of this formula. The boundary in Theorem~\ref{th:qiluo}, which answers
the source's direct question, is explicit.

The proof uses classical Laplace positivity and an elementary
hyperbolic inequality. The original source studies the fixed shift
$\mu=1/2$ with $\lambda\ge0$, and cites the classification without
the exponential correction. A bounded search did not locate the
full two-shift result. Its historical priority has not been established.



\section{Evidence and remaining questions}
For the two-data determinant, the earlier diffusion-wave result of
Floridia and Yamamoto, \emph{Backward problems in time for fractional
diffusion-wave equation},
\href{https://arxiv.org/abs/2007.09364}{arXiv:2007.09364}, concerns
the negative real ray. The rotating complex ray of
Theorem~\ref{th:wronskian} is the one in the original
Schr\"odinger problem. A search of cited and citing works and variants
of the determinant did not locate the same disproof for this ray.
The same authors still discuss the complete-monotonicity question in
their subsequent work \href{https://arxiv.org/abs/2603.16726}{arXiv:2603.16726}.
The contour expansion, winding number, and scale mixtures used here
have classical origins. The quantitative bounds and transfers are
stated separately so their novelty can be assessed directly.

The graph argument was checked separately against the degree-two chain
hypothesis in the Bankevich--Karpov theorem and against the small
component cases. The 2006 Foster abstract is a material gap in the
priority assessment. It discusses the same invariants but does not
identify all its results. A published source for the full linear bound
was not located.

For the normalized gamma-ratio derivative the argument uses
Kristiansen's theorem on gamma mixtures of shape two. Earlier work of
Hansen \cite{Hansen} also relates canonical convolution measures to
logarithmic curvature. The particular density comparison and the full
application were checked against the source formulas and against a
published statement giving an incompatible answer. The explicit gap
$r=2$ makes that disagreement verifiable independently. For
Qi--Wei--Guo, the work before the journal article dates to 2009. A
bounded search of citing works did not locate the sharp reciprocal
threshold. The Qi--Luo proof uses a classical Laplace representation
and a short hyperbolic inequality. These comparisons do not certify
a historically new general principle in every case.

The fractional-zero proofs combine analytical estimates with rigorously
rounded interval calculations. The exact statement with coefficient
one is distinct from the fitted constants in Conjectures 1 and 3.
Hanneken, Achar and Vaught \cite{HAV} had already computed a close
numerical value for the critical order. The results above certify its
location, establish positivity at all times below it and describe
both branches of the bifurcation locally. The implicit-function and
subordination inputs are classical. The uniform two-parameter estimate
is stated with additive error so its second-order terms can be tested.
Independent internal audits checked the finite and infinite portions
of the arguments. External peer review and exhaustive priority
verification are pending. 

All four interval-verification programs are embedded in the single TeX source
after \verb|\end{document}|. Copy each block delimited by its matching
\texttt{BEGIN} and \texttt{END} comments into its indicated Python file
and run it with Python-FLINT, which supplies Arb ball arithmetic. No
external figure, additional TeX file, or bibliography file is needed
for the PDF.



\appendix
\section{Typical shifts in the smooth-component question of Erd\H{o}s and Graham}
\label{sec:eg461-typical}

For an integer $t\ge 3$, let $s_t(m)=\prod_{p<t}p^{v_p(m)}$ for positive
integers $m$. Write $f(n,t)$ for the number of distinct $s_t(n+i)$ as
$1\le i\le t$. Erd\H{o}s and Graham's question (problem 461 in the
Erd\H{o}s Problems database, from \emph{Old and New Problems and Results
in Combinatorial Number Theory}, 1980, p.~92) asks for a positive lower
bound on $f(n,t)/t$ \emph{uniformly} in $n$ and $t$. We give an
asymptotic valid in mean square over the shifts. The uniform question
remains open in this work.

Let $H_j=\sum_{a=1}^j1/a$, $V_t=\prod_{p<t}(1-1/p)$, and
$M_t=\prod_{p<t}p^{\lceil\log t/\log p\rceil}$.
Let $\mathbb E_t$ denote the uniform average over a full period of
$M_t$ consecutive shifts.

\begin{theorem}\label{eg:typical}
As $t\to\infty$,
\[
 \mathbb E_t\left(\frac{f(n,t)}{t}-(1-V_tH_{t-1})\right)^2
   = O\left(\frac{\log\log t}{\log t}\right).
\]
Consequently $f(n,t)/t$ converges in mean square, and in natural-density
probability over the shifts, to $1-e^{-\gamma}$.
\end{theorem}

\begin{proof}
Every label $s_t(n+i)\ge t$ occurs once in the window, since equal
labels divide the difference of their positions. The events that the
label is $d$ for any $d<t$, or is at least $t$, are determined modulo
$M_t$. Indeed, reducing each valuation to the exponent in $M_t$ either
preserves the label or leaves a factor at least $t$. We may therefore
choose one independent uniform $p$-adic shift $N_p$ for each $p<t$.
The statistics used for the exact window count have precisely the
full-period law by the Chinese remainder theorem. Valuations are almost
surely finite. At one position let $S=\prod_{p<t}p^{v_p(N_p+i)}$.
For any $t$-smooth integer $d$ (all prime factors strictly below $t$),
independence and the geometric valuation law give
\[
     \Pr(S=d)=\frac{V_t}{d}.
\]
In particular, with $Z_i=\mathbf1_{\{S_i\ge t\}}$,
$\mathbb E Z_i=1-V_tH_{t-1}$.

Take $y=(\log t)^2$. At one position split $S=UX$ into its prime
powers for $p\le y$ and $y<p<t$, respectively. The factors $U,X$
are independent. Put $V_{y,t}=\prod_{y<p<t}(1-1/p)$. Conditional on
$U=m$, the chance that dropping $U$ changes the indicator of $S\ge t$
is at most
\[
 V_{y,t}\!\sum_{\substack{\max(1,t/m)\le d<t\\P^-(d)>y}}\frac1d,
\]
where $P^-(1)=\infty$. We use a harmonic packing bound. Multiply
each such $d$ by every integer $1\le b\le y$. The products $bd$
are distinct, since $b$ is their entire component supported on
primes up to $y$. All the products lie between $\max(1,t/m)$ and
$ty$. Comparing their harmonic sum with the harmonic sum over that
whole interval gives
\[
 H_{\lfloor y\rfloor}
 \sum_{\substack{\max(1,t/m)\le d<t\\P^-(d)>y}}\frac1d
 \le 1+\log m+\log y.
\]
The lost-indicator probability is at most
\[
  q_y=\frac{V_{y,t}}{H_{\lfloor y\rfloor}}
  \left(1+\log y+\sum_{p\le y}\frac{\log p}{p-1}\right)
  =(2+o(1))\frac{\log y}{\log t}.
\]
Here Mertens' two prime sum and product estimates apply whenever
$y=y(t)\to\infty$ and $y<t$. The estimate is independent of the
position.

For two positions distance $1\le h<t$ apart, compare their paired
vectors of valuations for $p>y$ with two independent vectors having
the correct marginals. If $p\nmid h$, simultaneous positive valuations
are impossible. Comparing the probabilities of the zero-zero cell,
the two axes, and the positive-positive cells gives total variation
distance $2/p^2$. If $p\mid h$, the distance is at most $4/p$ by
comparison to the point mass at $(0,0)$. The prime-wise pairs are
independent. Coupling each pair separately and using a union bound gives
total variation distance at most
\[
 2\sum_{p>y}p^{-2}+4\!\sum_{\substack{p>y\\p\mid h}}p^{-1}
   \ll \frac1y+\frac{\log t}{y\log y}.
\]
The last estimate follows since $h<t$ has at most $(\log t)/(\log y)$
distinct prime divisors larger than $y$. The covariance of the
indicators $\mathbf1_{\{X_i\ge t\}}$, $\mathbf1_{\{X_j\ge t\}}$
is bounded in absolute value by this total variation distance. Restoring
the small primes changes that covariance by at most $4q_y$. Consequently
\[
 \operatorname{Var}\left(t^{-1}\sum_{i=1}^t Z_i\right)
  \ll t^{-1}+q_y+y^{-1}+\frac{\log t}{y\log y}
  \ll \frac{\log\log t}{\log t}.
\]

It remains to count distinct labels smaller than $t$. Write $D$ for
their number and take $K=\lfloor t/\log t\rfloor$. There are at most
$K$ possible labels up to $K$. The number of distinct labels in
$(K,t)$ is at most the number of their occurrences, so
\[
   \frac{\mathbb ED}{t}\le \frac Kt+V_t\sum_{K<d<t}\frac1d
         \ll \frac{\log\log t}{\log t}.
\]
Since $0\le D\le t$, the same upper bound holds for
$\mathbb E[(D/t)^2]$. The exact identity $f=\sum_iZ_i+D$ and
$(a+b)^2\le2a^2+2b^2$ give the claimed estimate. Finally
$V_tH_{t-1}\to e^{-\gamma}$ by Mertens' product theorem.
\end{proof}

The publicly recorded averaging argument on the discussion page for
problem 461 obtains the upper bound $\inf_n f(n,t)/t\le
1-e^{-\gamma}+o(1)$. The pointwise marginal used above is also the
one relevant to that bound. For the original one-sided question a
concentration inequality gives a stronger exceptional-set estimate.

\begin{corollary}\label{eg:rare}
Fix $0<c<1-e^{-\gamma}$ and $0<a<(1-e^{-\gamma}-c)/2$. There is
$C_{a,c}>0$ such that for all sufficiently large $t$,
\[
 \Pr_t\{f(n,t)<ct\}\le \exp(-C_{a,c}t^a).
\]
The estimate holds after fixing arbitrary residue classes of the shift
at all prime powers in $M_t$ with base $p\le t^a$ and averaging only
over the remaining coordinates.
\end{corollary}
\begin{proof}
Use $y=t^a$. At each position the indicator of $X_i\ge t$ depends
only on the remaining prime coordinates. The number $B$ of positions
where this happens is a lower bound for $f$ for \emph{every} choice
of the small-prime residues. Indeed each such position must have
full smooth label at least $t$, and such labels are distinct. The
packing bound in the preceding proof yields
\[
   \mathbb EB/t\ \ge\ 1-V_tH_{t-1}-q_y
      \ \ge\ 1-e^{-\gamma}-2a-o(1)>c+\delta
\]
for some $\delta>0$. Changing the shift coordinate at one $p>y$
affects at most $2\lceil t/p\rceil$ entries, since the affected
positions are confined to the old and new zero residue classes
modulo $p$. Thus the sum of squares of coordinate oscillations of
$B/t$ is at most
\[
  \sum_{y<p<t} (2\lceil t/p\rceil/t)^2
       \le 16\sum_{m>y}m^{-2}\ll 1/y.
\]
The bounded-differences inequality for independent coordinates gives
$\Pr(B/t<c)\le\exp(-C_{a,c}y)$, as claimed. This inequality follows,
for example, by revealing the prime coordinates one at a time and
applying Hoeffding's exponential bound conditionally to the martingale
increments of $\mathbb E[B/t\mid\hbox{revealed coordinates}]$.
\end{proof}

\begin{corollary}\label{eg:almost-sure}
Let $c_0=1-e^{-\gamma}$ and fix $\epsilon>0$ and
$0<a<\min(1,\epsilon/3)$. Uniformly after fixing arbitrary residue
classes at all prime powers in $M_t$ with base $p\le t^a$, one has
\[
 \Pr_t\{|f(n,t)/t-c_0|>\epsilon\}
       \le 2\exp(-C_{a,\epsilon}t^a)
\]
for sufficiently large $t$. In particular, the periodic extension of $f$
to a Haar-random profinite integer $N$ satisfies $f(N,t)/t\to c_0$
almost surely through all integer $t$.
\end{corollary}
\begin{proof}
Keep $y=t^a$, with the coordinates at primes up to $y$ now fixed.
Let $U_i$ be the product of their prime powers at position $i$,
with each valuation capped at its exponent $e_p$ in $M_t$. If
$U_i<t$ no cap is effective. If $U_i\ge t$ the full label is
necessarily a singleton, independently of every higher valuation.
For these arbitrary fixed residues, counting multiples yields
\[
 \frac1t\sum_{i=1}^t\log\min(U_i,t)
 \le\sum_{p\le y}\log p\sum_{k=1}^{e_p}
                \Bigl(p^{-k}+t^{-1}\Bigr)
 \le \log y+O(1)+O(y\log t/t).
\]
Indeed $e_p\log p\le \log t+\log p\le2\log t$. To obtain a conditional
upper bound on the expected singleton count $\sum_i Z_i$, compare it
with $B$. At position $i$ a gain over $B$ requires
$t/\min(U_i,t)\le X_i<t$. The harmonic packing estimate with
$m=\min(U_i,t)$ bounds the probability by
\[
 \frac{V_{y,t}}{H_{\lfloor y\rfloor}}
      (1+\log y+\log\min(U_i,t)).
\]
Its average over $i$ is at most $2a+o(1)$ uniformly in the fixed
coordinates. Also $\mathbb EB/t\le 1-V_tH_{t-1}$, by first reintroducing
independent small-prime coordinates in the unrestricted calculation.
Consequently the conditional expected count $t^{-1}\sum_i Z_i$ is
between $c_0-2a-o(1)$ and $c_0+2a+o(1)$.

To bound the conditional expectation of $D/t$, again split at
$K=\lfloor t/\log t\rfloor$. If $U_i\ge t$ there is no small label
at this site. Otherwise the chance that the label belongs to $(K,t)$
is $V_{y,t}$ times a reciprocal sum over some of the $y$-rough integers
in $(K/U_i,t/U_i)$. Multiplying each term by all $b\le y$, as above,
bounds the latter sum by
\[
 \frac{1+\log(yt/K)}{H_{\lfloor y\rfloor}}.
\]
This works also when $K/U_i<1$, since the lower endpoint of the
harmonic sum after multiplication is then at least one. Adding $K/t$
gives $\mathbb E[D/t\mid\text{fixed coordinates}]\le a+o(1)$.
Combining both bounds places the conditional expectation of $f/t$
between $c_0-2a-o(1)$ and $c_0+3a+o(1)$, uniformly for the
specified small-prime residues.

Changing the coordinate of one prime $p>y$ alters at most
$2\lceil t/p\rceil$ sites. Changing $m$ site labels in a multiset
alters its number of distinct labels by at most $m$. One may use the
truncated labels $\gcd(n+i,M_t)$ directly for this argument, since
all labels at least $t$ are singletons even after truncation. Thus the
bounded-differences estimate from Corollary~\ref{eg:rare} applies to
$f/t$ itself. This gives the asserted two-sided bound because
$\epsilon>3a$. Summing the resulting unrestricted upper bound over
$t$ and invoking the first Borel--Cantelli lemma for each positive
rational $\epsilon$ proves the almost-sure claim.
\end{proof}

For comparison, each fixed ordinary shift $n\ge0$ satisfies
$f(n,t)/t\to1$, since the $m<t$ in its window all equal their own
smooth components. At $n=0$ the number of labels at least $t$ is at
most one, although there are at least $t-1$ distinct labels. Thus the
singleton estimate used in the concentration proof cannot give a
uniform linear lower bound by itself. The natural density and fixed-shift limits have
different quantifiers.

An exceptional set can have positive cardinality even when its
proportion of the full period is exponentially small in a power of
$t$. The estimate leaves the uniform lower-bound question unresolved.

\section{Free isospectral potentials with four occupied sites}
\label{sec:four-sites}
Faust, Friedman, O'Malley, Ramos, and Sharma ask in Conjecture~2 of
\emph{Algebraic properties of the ideal of spectral invariants for the
discrete Laplacian} (arXiv:2602.08304) for a nonzero complex
$q$-periodic potential isospectral to the free one and satisfying a
conjugate-reflection symmetry, for every $q>3$.
Faust et al., J. Math. Phys. \textbf{65} (2024), 073501,
constructed such even-period potentials using four occupied sites.
We determine exactly when a four-site construction is possible,
allowing arbitrary complex values. The existence question for odd
periods and unrestricted support remains open in this work.

For $t\ne0$ and $x=t+t^{-1}$ write
$[j]=(t^j-t^{-j})/(t-t^{-1})$. For positive $j$ this is a monic
polynomial of degree $j-1$ in $x$. The two-sided recurrence is
$[j+1]=x[j]-[j-1]$, with $[0]=0$ and $[1]=1$.
The discriminant is the trace of the ordered product of the matrices
$T_j=\left(\begin{smallmatrix}x-v_j&-1\\1&0\end{smallmatrix}\right)$.
Equality of the discriminants is equivalent to Floquet isospectrality
since the transfer matrices have determinant one.

\begin{theorem}
Let $q\ge4$. A nonzero free-isospectral complex potential supported on at
most four sites has exactly four occupied sites. Up to cyclic choice of
the initial occupied site, their consecutive values must be
\[
 (A,-A,B,-B),\qquad B=i\epsilon A,\quad\epsilon\in\{-1,1\}.
\]
Write $(a,b,c,d)$ for the consecutive positive cyclic gaps, starting with
the edge from $A$ to $-A$. Then $b=d=1$, and exactly the following
possibilities occur:
\begin{enumerate}
\item $a=c\ge1$, $q=2a+2$, $A^2=2i\epsilon$;
\item $\{a,c\}=\{1,2\}$, $q=5$,
$A^2=2i\epsilon+\operatorname{sgn}(c-a)$.
\end{enumerate}
Conversely every listed pattern is free-isospectral. The second case
does not have a conjugate-invariant multiset of entries. In particular,
for every odd $q\ge5$ any potential satisfying the original conjecture
requires at least five occupied sites.
\end{theorem}
\begin{proof}
Expand the discriminant multilinearly using $T_j=T-v_je_1e_1^{\mathsf T}$,
where $T$ is the free transfer matrix. For any nonempty set $I$ of $k$
occupied sites, with successive cyclic gaps $g_1,\ldots,g_k$, the contribution
to the difference from the free discriminant is
\[
 (-1)^k\prod_{j\in I}v_j\,\prod_{l=1}^k[g_l].
\]
Indeed $e_1^{\mathsf T}T^{g-1}e_1=[g]$, as follows from the recurrence.
The product of gap polynomials has degree $q-k$, parity $q-k$, and
leading coefficient one. The $k=1$ group is $-[q]\sum v_j$, so the
highest coefficient forces $e_1(v)=0$ and this group vanishes identically.
The next two coefficients therefore force $e_2(v)=e_3(v)=0$ (at the
third coefficient the $k=2$ group contributes nothing by parity).
With at most three occupied sites, these identities force every entry
to be zero. With four occupied sites, the four values are the roots of
$X^4-A^4$ for some $A\ne0$, namely $A,iA,-A,-iA$.

The $k=3$ group has opposite parity to the groups $k=2,4$, hence must
vanish individually. Denote the four cyclic values by $v_j$, and the
edge lengths from $j$ to $j+1$ by $g_j$. Dividing the triple sum by
$\prod_{j=1}^4 v_j[g_j]$ gives
\[
 \sum_{j=1}^4 v_j^{-1}
    \frac{[g_{j-1}+g_j]}{[g_{j-1}][g_j]}=0.
\]
The ratio $[u+v]/([u][v])$ equals
$(t-t^{-1})(H_u+H_v)/2$, where $H_u=(t^{2u}+1)/(t^{2u}-1)$.
The rational functions $H_u$ for distinct positive $u$ are linearly
independent: a primitive $2u$-th root is a pole of the largest-index
term and no smaller one. Iteration proves independence. Consequently
for every occurring length $g$,
\[
 \sum_{j:g_j=g}(v_j^{-1}+v_{j+1}^{-1})=0.                 \tag{*}
\]

If no cyclically neighboring values are opposites, we may list the
values as $(A,B,-A,-B)$ with $B=i\epsilon A$. In (*) the four edge
weights consist of two pairs of opposites, represented by
$A^{-1}+B^{-1}$ and $B^{-1}-A^{-1}$, which are linearly independent
over the reals. Thus the cyclic gaps are $(a,b,a,b)$.
The residual from the pair and quartet groups is
\[
 2i\epsilon A^2\bigl([a][q-a]-[b][q-b]\bigr)
                  -A^4[a]^2[b]^2.
\]
Here $q=2a+2b$ and the difference in parentheses equals
$[a-b][a+b]=[a]^2-[b]^2$. Its vanishing would require
$A^2=2i\epsilon([b]^{-2}-[a]^{-2})$ as a rational-function identity.
If both $a,b>1$ the right side tends to zero at infinity. If exactly
one is one it varies with $x$, while if both are one it vanishes.
This rules out this cyclic order.

Otherwise the cyclic values can be written $(A,-A,B,-B)$ with
$B=i\epsilon A$. Formula (*) gives gaps $(a,b,c,b)$. Directly summing
the six pair terms and the quartet now gives, after division by
$A^2[a][c]$,
\[
 \frac{[c-a][2b]}{[a][c]}+2i\epsilon-A^2[b]^2=0.       \tag{**}
\]
To verify this simplification, the opposite-pair terms before division
are $A^2([c][a+2b]-[a][c+2b])=A^2[c-a][2b]$, and
$[a+b][b+c]-[b][a+b+c]=[a][c]$ for the two cross-pair lengths.
These identities follow by expanding brackets in $t$.
For real $x>2$, take imaginary parts of (**). Since $2\epsilon\ne0$,
$[b]^2$ must be constant, giving $b=1$.
If $a=c$ this gives the first family. For $a\ne c$ the nonzero
ratio $[c-a][2]/([a][c])$ has degree $2-2\min(a,c)$ at infinity,
so it can only be constant if $\min(a,c)=1$. With $u=\max(a,c)>1$
its absolute sign-adjusted expression is $x[u-1]/[u]$. It tends to
one at infinity, and the recurrence $x[u-1]=[u]+[u-2]$ shows it
is identically one exactly when $u=2$. This gives the second family.
All steps after the vanishing of $e_1,e_2,e_3$ characterize the
discriminant identity, proving the converse. In the second case
$A^4=-3+4i\epsilon\operatorname{sgn}(c-a)$ is nonreal, so the
polynomial $X^4-A^4$ is not real. Thus conjugation does not preserve
the multiset of its roots.
\end{proof}

\paragraph{An independent support-width obstruction.}
If the shortest cyclic interval covering the support of a nonzero
$q$-periodic free-isospectral potential contains $h$ sites, then
$h\ge\lfloor q/2\rfloor+2$. The even four-site family above attains
this bound. To see this, place the occupied endpoints of the interval
at $1,h$, write $M=T_h\cdots T_1$, and use the eigenbasis
$S=\left(\begin{smallmatrix}t&t^{-1}\\1&1\end{smallmatrix}\right)$
of $T$. Set $U(t)=(S^{-1}MS)_{11}$ and
$P(t)=(t-t^{-1})(U(t)-t^h)$. The other diagonal entry is $U(1/t)$.
Equality of the two period discriminants thus means
$t^{q-h}P(t)=t^{h-q}P(t^{-1})$. On expanding the transfer product,
the contribution of $j_1<\dots<j_k$ to $P(t)$ is
\[
 (-1)^k\prod v_{j_i}\,t^{h-j_k+j_1}
 \prod_{i=1}^{k-1}[j_{i+1}-j_i].
\]
Its smallest exponent equals $h-2(j_k-j_1)+k-1$. For $h\ge2$ the
unique smallest exponent over all subsets is $3-h$, achieved only
by the pair $\{1,h\}$, with nonzero coefficient $v_1v_h$.
Therefore the Laurent polynomial $t^{q-h}P(t)$ has lowest exponent
$q-2h+3$. Since it equals its reciprocal and is nonzero,
$q-2h+3\le0$. This gives the stated bound. When $h=1$ the
discriminant difference is simply $-v_1[q]$.

The 2024 construction already gives the even-period family and the
2026 source lists computed solutions for period five. Corpus [061]
also gives many individual odd periods. The all-period four-site
classification and width bound above address support restrictions.
The original existence question at arbitrary odd periods with
unrestricted support remains open in this work.



\begin{thebibliography}{99}
\bibitem{wow165:registry}
E. DeLaVi\~na, \emph{Written on the Wall II}, Graffiti.pc
Conjecture~165, 8 August 2005.
\href{http://cms.dt.uh.edu/faculty/delavinae/research/wowII/}{Original catalogue}.
An archived copy of the HTML row with its Foster annotation is available
in the \href{https://gitlab.com/ai-village-agents/village/graffiti-verification/-/blob/3be380b4222f974bedbefb228c8e75fbd9a7b3ae/data/wow2_rows_html.json}{Graffiti verification archive}.

\bibitem{wow165:BK}
A. Bankevich and D. Karpov,
\emph{Bounds of a number of leaves of spanning trees},
arXiv:1111.3266 (2011), Theorem~2.
\url{https://arxiv.org/abs/1111.3266}.

\bibitem{wow165:Foster}
T. C. Foster, \emph{On Leafy Spanning Trees of Graph},
5th Annual Student Research Conference, University of Houston--Downtown
(2006), conference abstract.
\url{https://www.uhd.edu/documents/academics/sciences/scholars/proceedings2006.pdf}.

\bibitem{CDH}
R. Chaudhary, K. Diethelm and S. Hashemishahraki,
\emph{On the separation of solutions to fractional differential equations
of order $\alpha\in(1,2)$}, Applied Numerical Mathematics 203 (2024), 84--96.
\href{https://doi.org/10.1016/j.apnum.2024.05.020}{doi:10.1016/j.apnum.2024.05.020}.
\href{https://arxiv.org/abs/2401.14771}{arXiv:2401.14771}.
\bibitem{GM}
R. Gorenflo and F. Mainardi,
\emph{Fractional calculus: integral and differential equations of
fractional order}, in Fractals and Fractional Calculus in Continuum
Mechanics (1997), 223--276. \url{https://arxiv.org/abs/0805.3823}.
\bibitem{HAV}
J. W. Hanneken, B. N. N. Achar and D. M. Vaught, \emph{An alpha-beta phase diagram representation of the zeros and properties
of the Mittag--Leffler function}, Advances in Mathematical Physics
(2013), article 421685.
\href{https://doi.org/10.1155/2013/421685}{doi:10.1155/2013/421685}.
\bibitem{Hansen}
B. G. Hansen, \emph{On log-concave and log-convex infinitely divisible
sequences and densities}, Annals of Probability 16 (1988), 1832--1839.
\href{https://doi.org/10.1214/aop/1176991600}{doi:10.1214/aop/1176991600}.
\bibitem{EG}
P. Erd\H{o}s and R. L. Graham,
\emph{Old and New Problems and Results in Combinatorial Number Theory},
Monographies de l'Enseignement Math\'ematique 28 (1980), p.~92.
\url{https://www.erdosproblems.com/461}.
\end{thebibliography}
\end{document}

%% BEGIN schrodinger1.py
"""An interval certificate of a purely imaginary Mittag-Leffler zero.
Run with python-flint 0.9: only rounded-out Arb balls and exact integers."""
from math import factorial
from flint import arb, acb, ctx

ctx.prec = 160
radius = arb(10) ** -12
alpha0 = arb("0.7478564279599293430791082426")
t0 = arb("3.7131384963793705236713833185")
box = arb(0, radius)
tail = arb(10) ** -42
# Pair k=2m,2m+1; Gamma(1+alpha*k) >= m!, k>=240, alpha>1/2.
# Each invisible term incl. first derivatives has size <= (k+1)^2*4^k/m!.
# Ratio of successive paired majorants <= (16/121)*(244/242)^2 < 1/4.
assert 4*16*244**2 < 121*242**2
assert 7 * 242**2 * 4**240 * 10**42 < factorial(120)

def evaluate(alpha, t):
    power = acb(1)
    z = acb(0, t)
    output = [acb(0)] * 3
    for k in range(240):
        u = 1+k*alpha
        v = u.rgamma()*power
        output[0] += v
        output[1] -= k*u.digamma()*v
        output[2] += k*v/t
        power *= z
    return [(v.real + arb(0,tail), v.imag + arb(0,tail)) for v in output]

center = evaluate(alpha0,t0)[0]
df_a,df_t = evaluate(alpha0+box,t0+box)[1:]
B = ((arb(424)/1000,-arb(523)/1000), (arb(8133)/1000,arb(792)/1000))
# T = x-B F(x). Bound infinity norms of T(x0)-x0 and T' over entire box.
shift=max(abs(b[0]*center[0]+b[1]*center[1]) for b in B)
q=arb(0)
for i, row in enumerate(B):
    norm=arb(0)
    for j, col in enumerate((df_a,df_t)):
        norm+=abs(int(i==j)-row[0]*col[0]-row[1]*col[1])
    q=max(q,norm)
print("shift:",shift,"contraction:",q)
assert shift < radius/1000 and q < arb(1)/100

%% END schrodinger1.py

%% BEGIN schrodinger2.py
"""CETMY Conjecture 2. Certified zero of its original complex determinant."""
from flint import arb, acb, ctx
from math import factorial
ctx.prec = 176
rad = arb(10)**-12
ac = arb("1.68430936268566911574800653125")
wc = arb("5.40314724546132388337747818913")
pad = arb(0,rad)
err = arb(0,arb(10)**-36)
N = 128
# For k>=N, Gamma(a*k+b)>=k!, |z|<20, and the factors for both
# first derivatives are <= 3*(k+1)^2. Majorant ratio is < 1/2.
assert 6 * (N+1)**2 * 20**N * 10**36 < factorial(N)
assert 40*(N+2)**2 < (N+1)**3
def det3(a,w):
    L=acb(w.log(),-arb.pi()/2)
    z=(a*L).exp()
    assert abs(z)<20 and abs(L)<4 and a>1 and a<2 and w>5
    vals=[]
    for b,db in ((arb(1),0),(arb(2),0),(a,1)):
        f=fa=fw=acb(0)
        power=acb(1)
        for k in range(N):
            u=a*k+b
            term=power*u.rgamma()
            f+=term
            fa+=term*(k*L-(k+db)*u.digamma())
            fw+=term*(a*k/w)
            power*=z
        vals.append([h+acb(err,err) for h in (f,fa,fw)])
    (f,fa,fw),(g,ga,gw),(h,ha,hw)=vals
    D=f*f-z*g*h
    Da=2*f*fa-z*(L*g*h+ga*h+g*ha)
    Dw=2*f*fw-z*(a*g*h/w+gw*h+g*hw)
    return D,Da,Dw
center=det3(ac,wc)[0]
da,dw=det3(ac+pad,wc+pad)[1:]
J=((da.real,dw.real),(da.imag,dw.imag))
B=((arb(896)/1000,-arb(459)/1000),(arb(10189)/1000,arb(14195)/1000))
shift=max(abs(b[0]*center.real+b[1]*center.imag) for b in B)
q=max(sum(abs(int(i==j)-sum(B[i][k]*J[k][j] for k in range(2))) for j in range(2)) for i in range(2))
assert abs(B[0][0]*B[1][1]-B[0][1]*B[1][0])>0
print("shift",shift,"contraction",q)
assert shift<rad/1000 and q<arb(1)/100

%% END schrodinger2.py

%% BEGIN cdh_cert.py
"""Uniform first-zero bounds, evaluated in outward-rounded Arb balls."""
from flint import arb, ctx
from math import factorial
ctx.prec=112
N=60
tail=arb(10)**-13
q=arb(9)**(N+1)/factorial(N+1)
# Absolute errors for F, F_X and d/dalpha F(alpha,S**alpha)
assert 8*(N+2)**2*q < tail

def sphere(a,b):
    return a.union(b)

def horner(coeff,x):
    r=arb(0)
    for c in reversed(coeff):
        r=r*x+c
    return r

def check(amin,amax,mode,S):
    # Monotone separate digamma and reciprocal gamma endpoints.
    A=sphere(amin,amax); B=A if mode else arb(1)
    mid=(amin+amax)/2
    X=(A*S.log()).exp(); Xm=(mid*S.log()).exp()
    c=[]; cm=[]; dp=[]
    for k in range(N+1):
        l=amin*(k+mode)+(1-mode); r=amax*(k+mode)+(1-mode)
        if k==0 and mode:
            cc=A.rgamma()
        else:
            cc=sphere(l.rgamma(),r.rgamma()) if l!=r else l.rgamma()
        c.append((-1)**k*cc)
        cm.append((-1)**k*(mid*(k+mode)+(1-mode)).rgamma())
        psi=sphere(l.digamma(),r.digamma()) if l!=r else l.digamma()
        dp.append(c[-1]*(k*S.log()-(k+mode)*psi))
    # Centered alpha evaluation at the right edge of the t interval.
    val=horner(cm,Xm)+sphere(amin-mid,amax-mid)*horner(dp,X)
    if not (val-tail*(1+amax-amin)>0):return False
    deriv=[(k+1)*c[k+1] for k in range(N)]
    # The x derivative is negative from zero to S**alpha.
    maxx=(amax*S.log()).exp()
    todo=[(arb(0),maxx,0)]
    while todo:
        a,b,d=todo.pop()
        bound=horner(deriv,sphere(a,b))+tail
        if bound<0:continue
        if d==8:return False
        m=(a+b)/2
        todo.extend([(a,m,d+1),(m,b,d+1)])
    return True

if __name__=='__main__':
    for mode,S in ((1,arb(293785)/100000),(0,arb(1559)/1000)):
        passed=0; depth=0; todo=[(arb(1),arb(2),0)]
        while todo:
            a,b,d=todo.pop()
            if check(a,b,mode,S):
                passed+=1;depth=max(depth,d);continue
            if d>=22:raise ArithmeticError((mode,a,b))
            m=(a+b)/2;todo.extend([(a,m,d+1),(m,b,d+1)])
        print('beta_alpha',mode,'passed boxes',passed,'max depth',depth,flush=True)

if __name__=='__main__':
    # At alpha=2 the two zeros are simple and F_t is known exactly.
    for mode,low,high in ((1,-810252,-810251),(0,-129292,-129291)):
        t=arb.pi() if mode else arb.pi()/2
        b=1+mode
        derivative=arb(0)
        for k in range(60):
            s=arb(2*k+b)
            derivative+=(-1)**k*t**(2*k)*s.rgamma()*(k*t.log()-(k+mode)*s.digamma())
        rem=arb(12)*61**2*16**60/factorial(120)
        assert rem<arb(10)**-65
        slope=-(t if mode else 1)*(derivative+arb(0,rem))
        assert slope>arb(low)/10**6 and slope<arb(high)/10**6
        print('beta_alpha',mode,'slope',slope)

%% END cdh_cert.py

%% BEGIN cdh_fold.py
"""Interval checks for the global fold of E_{a,2}(-t**a). Python-FLINT."""
from flint import arb, ctx
from fractions import Fraction as Q
from functools import lru_cache
ctx.prec=144
N=70
def R(x):
    q=Q(x)
    return arb(q.numerator)/q.denominator
def hull(l,u):return R(l).union(R(u))
def hor(cs,x):
    y=arb(0)
    for c in reversed(cs):y=y*x+c
    return y
def err(v):return v+arb(0,arb(10)**-36)

# Tests use a>=1.59, t<=10, and reciprocals t^{-1} only for t>=1.
# The kth differentiated term in this region is bounded by
# 20(k+2)^3 *100^k/Gamma(1.59*k+2) IF t^a<=100.
# Tighter for the tail, consider w=t^(a*k)/Gamma(a*k+2):
# gamma ratio gives w(k+1)/w(k)<= (10/(a*k+2))^a.
# For a*k>10, w decreases with a since psi(a*k+2)>log(10).
start=arb(N+1)*R('1.59')+2
assert start.digamma()>arb(10).log()
assert (arb(10)/(start))**R('1.59')*(R(N+4)/R(N+3))**3<arb(1)/2
tail=40*(N+3)**3*(R('1.59')*(N+1)*arb(10).log()).exp()*start.rgamma()
assert tail<R('1e-36')

@lru_cache(maxsize=None)
def coeff(alo,ahi):
    result=[]; psis=[]
    for k in range(N+1):
        l=R(alo)*k+2;h=R(ahi)*k+2
        c=l.rgamma().union(h.rgamma()) if alo!=ahi else l.rgamma()
        psi=l.digamma().union(h.digamma()) if alo!=ahi else l.digamma()
        result.append(c*(-1)**k); psis.append(psi)
    return result,psis

def evalbox(al,au,tl,tu,need):
    a=hull(al,au); t=hull(tl,tu); logt=t.log(); x=(a*logt).exp()
    c,psi=coeff(al,au)
    out={}
    for key in need:
        if key=='F':out[key]=err(hor(c,x));continue
        if key=='Ft':
            s=[a*k*v for k,v in enumerate(c)];out[key]=err(hor(s,x)/t);continue
        if key=='Ftt':
            s=[a*k*(a*k-1)*v for k,v in enumerate(c)];out[key]=err(hor(s,x)/t**2);continue
        if key=='Fa':
            s=[k*(logt-psi[k])*v for k,v in enumerate(c)];out[key]=err(hor(s,x));continue
        if key=='Fat':
            s=[(1+a*k*(logt-psi[k]))*k*v for k,v in enumerate(c)];out[key]=err(hor(s,x)/t);continue
    return out

ac,tc=Q('1.599115206323013'),Q('5.210687979825734')
ea,et=Q('1e-13'),Q('1e-11')
al,au=ac-ea,ac+ea
tl,tu=tc-et,tc+et
C=evalbox(ac,ac,tc,tc,['F','Ft','Fa','Ftt','Fat'])
J=evalbox(al,au,tl,tu,['Fa','Ft','Fat','Ftt'])
det=C['Fa']*C['Ftt']-C['Ft']*C['Fat']
Y=[[ (C['Ftt']/det).mid(),(-C['Ft']/det).mid()],
   [(-C['Fat']/det).mid(),(C['Fa']/det).mid()]]
g=[C['F'],C['Ft']]; M=[[J['Fa'],J['Ft']],[J['Fat'],J['Ftt']]]
rad=[R(ea),R(et)]
for i in range(2):
    disp=abs(sum(Y[i][j]*g[j] for j in range(2)))
    for j in range(2):
        disp+=abs(int(i==j)-sum(Y[i][k]*M[k][j] for k in range(2)))*rad[j]
    assert disp<rad[i]
assert abs(Y[0][0]*Y[1][1]-Y[0][1]*Y[1][0])>0
cusp=-( -2*J['Fa']/J['Ftt']).sqrt()
assert R('-4.789')<cusp and cusp<R('-4.787')
print('fold box and sqrt coefficient:',cusp,flush=True)

def verify_1d(label,l,u,pred,maxdepth=24):
    todo=[(l,u,0)];cnt=0;depth=0
    while todo:
        x,y,d=todo.pop()
        if pred(x,y):cnt+=1;depth=max(d,depth);continue
        if d==maxdepth: raise ArithmeticError((label,x,y))
        mid=(x+y)/2
        todo.extend([(x,mid,d+1),(mid,y,d+1)])
    print(label,cnt,depth,flush=True)

# At critical a, the only zero on 3<=t<=10 is in the fold neighborhood.
def poscrit(l,u):
    mid=(l+u)/2
    D=evalbox(al,au,l,u,['Ft'])['Ft']
    F=evalbox(al,au,mid,mid,['F'])['F']
    return F-abs(D)*R((u-l)/2)>0
verify_1d('critical positivity left',Q(3),Q('5.20'),poscrit)
verify_1d('critical positivity right',Q('5.22'),Q(10),poscrit)
verify_1d('critical convex strip',Q('5.20'),Q('5.22'),
          lambda l,u:evalbox(al,au,l,u,['Ftt'])['Ftt']>0)

# For all later orders, F is positive on t<=3. Use X=t^a.
def pos_upto_three(l,u):
    c,_=coeff(l,u)
    xmax=(R(u)*R(3).log()).exp()
    # Value positive through t=3 unless a slightly exceeds 2 (not used).
    X=(hull(l,u)*R(3).log()).exp()
    if not (err(hor(c,X))>0):return False
    cs=[(j+1)*c[j+1] for j in range(N)]
    todo=[(arb(0),xmax,0)]
    while todo:
        x,y,d=todo.pop()
        if err(hor(cs,x.union(y)))<0:continue
        if d==6:return False
        m=(x+y)/2;todo.extend([(x,m,d+1),(m,y,d+1)])
    return True
verify_1d('all later orders, t<=3',Q('1.5991'),Q(2),pos_upto_three)

# Strict decrease in the order across the strip containing all first zeros.
def neg_da(l,u):
    todo=[(Q(3),Q('5.22'),0)]
    while todo:
        x,y,d=todo.pop()
        if evalbox(l,u,x,y,['Fa'])['Fa']<0:continue
        if d==9:return False
        m=(x+y)/2;todo.extend([(x,m,d+1),(m,y,d+1)])
    return True
verify_1d('order derivative negative',Q('1.5991'),Q(2),neg_da)

a=hull(al,au); phase=arb.pi()/a
A=-(arb.pi()*a).sin()/(arb.pi()*(a-1)*arb(1).exp()*(1+arb(1)/10)**2)
assert A*arb(10)**(1-a)>(2/a)*(10*phase.cos()).exp()
assert phase.cos()+(a-1)/10<0
print('infinite tail bound passed',flush=True)

%% END cdh_fold.py
